% BEGIN EMBEDDED BIBLIOGRAPHY
% Bibliography data are part of this source; the compiler writes its temporary input.
\begin{filecontents*}[overwrite]{\jobname.bib}
@unpublished{QED3Master,
 title={The field-strength {Osterwalder--Schrader} axioms for {QED} in three dimensions},
 shorthand={QED}, date={2026}, note = {Comprehensive edition, 263 pp., 20 October 2026, before the modular integration revision}}
@unpublished{MaxwellBlockManuscript,
 title={Maxwell block minimization and Gaussian fluctuations},
 shorthand={MB}, date={2026}, note={Companion manuscript, version of 20 October 2026}}
@unpublished{FermionBlockManuscript,
 title={Chamber fermion locality and local integration},
 shorthand={FB}, date={2026}, note={Companion manuscript, version of 20 October 2026}}
@unpublished{NormalizedContraction,
 title={Relevant extraction and contraction of local interactions},
 shorthand={NC}, date={2026}, note={Companion manuscript, version of 20 October 2026}}
@misc{PostnikovLin2019,
 author={Postnikov, Alexander}, title={Algebraic Combinatorics},
 date={2019-04-05}, note={MIT 18.212, Lecture 23, notes by Andrew Lin, Proposition 4},
 url={https://ocw.mit.edu/courses/18-212-algebraic-combinatorics-spring-2019/f6113618a3e085ba86b799f65934886d_MIT18_212S19_lec23.pdf}}
\end{filecontents*}
% END EMBEDDED BIBLIOGRAPHY
% NLM1 / d040: full proof companion to NLO1.
\documentclass[11pt,letterpaper,reqno]{amsart}
% BEGIN INLINE manuscript-style
\makeatletter
\NeedsTeXFormat{LaTeX2e}[2020/10/01]
% Use with \documentclass[11pt,letterpaper,reqno]{amsart}.
% This file controls presentation only. It supplies no mathematical content.

\RequirePackage[T1]{fontenc}
\RequirePackage{lmodern}
\RequirePackage{microtype}
\RequirePackage[letterpaper,margin=1in,headheight=15pt]{geometry}
\RequirePackage{amsmath,amsthm,amssymb,mathtools}
\RequirePackage{enumitem}
\RequirePackage{perpage}
\RequirePackage{xcolor}
\definecolor{manuscriptred}{rgb}{0.6,0.1,0.1}
\definecolor{manuscriptblue}{rgb}{0.1,0.3,0.5}
\definecolor{manuscriptmagenta}{HTML}{60445C}
\RequirePackage[breaklinks=true,colorlinks=true,
  linkcolor=manuscriptred,citecolor=manuscriptblue,urlcolor=manuscriptmagenta]{hyperref}
\RequirePackage{cleveref}
\hypersetup{pdfauthor={}}

% Retain the AMS paragraph, title, heading and running-head conventions.
% Equations retain one global Arabic sequence, as in the reference.
\setcounter{tocdepth}{1}
\renewcommand{\thefigure}{\Roman{figure}}
\renewcommand{\thetable}{\Roman{table}}
\MakePerPage{footnote}
\renewcommand{\thefootnote}{%
  \ifnum\value{footnote}>9\arabic{footnote}\else\fnsymbol{footnote}\fi}

% Upright statements. Main results are A, B, C; supporting statements
% share a section-local sequence. These end marks indicate typography,
% never proof completeness, validation, or a deferred-proof status.
\theoremstyle{definition}
\newtheorem{ms@theorem}{Theorem}
\renewcommand{\thems@theorem}{\Alph{ms@theorem}}
\newtheorem*{ms@theoremstar}{Theorem}
\newtheorem{ms@proposition}{Proposition}[section]
\newtheorem{ms@lemma}[ms@proposition]{Lemma}
\newtheorem{ms@definition}[ms@proposition]{Definition}
\newtheorem{ms@corollary}[ms@proposition]{Corollary}
\newtheorem*{ms@corollarystar}{Corollary}

\newcommand{\ms@majorstart}{%
  \pushQED{\qed}\renewcommand{\qedsymbol}{$\blacksquare\square$}}
\newcommand{\ms@minorstart}{%
  \pushQED{\qed}\renewcommand{\qedsymbol}{$\blacksquare$}}
\newenvironment{theorem}[1][]{\ms@majorstart\ms@theorem[#1]}
  {\popQED\endms@theorem}
\newenvironment{theorem*}[1][]{\ms@majorstart\ms@theoremstar[#1]}
  {\popQED\endms@theoremstar}
\newenvironment{proposition}[1][]{\ms@minorstart\ms@proposition[#1]}
  {\popQED\endms@proposition}
\newenvironment{lemma}[1][]{\ms@minorstart\ms@lemma[#1]}
  {\popQED\endms@lemma}
\newenvironment{definition}[1][]{\ms@minorstart\ms@definition[#1]}
  {\popQED\endms@definition}
\newenvironment{corollary}[1][]{\ms@majorstart\ms@corollary[#1]}
  {\popQED\endms@corollary}
\newenvironment{corollary*}[1][]{\ms@majorstart\ms@corollarystar[#1]}
  {\popQED\endms@corollarystar}
% The source's lemma/proposition variants are deliberately not carried over:
% choose one marker for each statement type rather than encode proof status.

\theoremstyle{remark}
\newtheorem{remark}[ms@proposition]{Remark}
\newtheorem*{remark*}{Remark}
% The ordinary amsthm proof environment retains its open square.

\crefname{ms@theorem}{theorem}{theorems}
\Crefname{ms@theorem}{Theorem}{Theorems}
\crefname{ms@proposition}{proposition}{propositions}
\Crefname{ms@proposition}{Proposition}{Propositions}
\crefname{ms@lemma}{lemma}{lemmas}
\Crefname{ms@lemma}{Lemma}{Lemmas}
\crefname{ms@definition}{definition}{definitions}
\Crefname{ms@definition}{Definition}{Definitions}
\crefname{ms@corollary}{corollary}{corollaries}
\Crefname{ms@corollary}{Corollary}{Corollaries}
\crefname{remark}{remark}{remarks}
\Crefname{remark}{Remark}{Remarks}

% Bibliography configuration for this manuscript.
  \RequirePackage[backend=biber,style=alphabetic,backref=true,
    giveninits=true,isbn=false,date=year]{biblatex}
  \addbibresource{\jobname.bib}
  \setcounter{biburlnumpenalty}{100}
  \setcounter{biburllcpenalty}{100}
  \setcounter{biburlucpenalty}{100}
  \renewcommand*{\bibfont}{\footnotesize}
  \renewbibmacro{in:}{}
  \DeclareFieldFormat{pages}{#1}
  % Use the public alphabetic style. No private author-hash overrides;
  % repeated author names are printed in full.
\makeatother
% END INLINE manuscript-style
\usepackage{mathrsfs}
\newcommand{\R}{\mathbb R}
\newcommand{\C}{\mathbb C}
\newcommand{\Z}{\mathbb Z}
\newcommand{\N}{\mathbb N}
\newcommand{\Sspace}{\mathcal S}
\newcommand{\dd}{\mathrm d}
\DeclareMathOperator{\supp}{supp}
\DeclareMathOperator{\tr}{tr}
\DeclareMathOperator{\diag}{diag}
\newcommand{\extref}[2]{#1 of~\cite{#2}}
\newcommand{\exteq}[2]{(#1) of~\cite{#2}}
\allowdisplaybreaks
\title[The exact nonlinear renormalization step]
{The exact nonlinear renormalization step\\for lattice QED in three dimensions}
\date{}
\hypersetup{pdftitle={The exact nonlinear renormalization step for lattice QED in three dimensions},pdfauthor={}}
\begin{document}
\begin{abstract}
We prove an exact finite-volume block identity for lattice QED with
chamber fermions and an independent external connection. Under explicit
domain, primitive and small-charge hypotheses, two finite localizations
produce a fluctuation activity whose connected logarithm converges
absolutely. We prove its local symmetries, a norm bound after scaling,
and a coupling-difference estimate on the unit data ball. The resulting
activity, mass, vacuum and external scalar maps reproduce the complete
outgoing density and its unnormalized photon-box marginal. The first
step integrates all chamber components with their weighted Berezin
measure.
\end{abstract}
\maketitle
\tableofcontents
\section{Introduction}\label{sec:introduction}

We consider one block transformation of noncompact lattice QED with
four-component chamber fermions. The photon field is periodic on
$\eta\Z^3/\ell\Z^3$, and spinors are antiperiodic in every coordinate.
Each photon site carries $48$ chamber labels. The chamber measure is
$\eta^3/48$; barred and unbarred Grassmann generators are independent.
The Dirac stencil, block averages and weighted transposes are those of
\cite[\S1]{FermionBlockManuscript}.

Fix \(0<\rho<1\), odd \(L\ge4^8\), an integer power \(M\) of \(L\) with \(M\ge L^2\) and \(M>4L\),
and \(k\in\Z_{\ge0}\). Put \(\eta=L^{-k}\), \(e=e_k\in(0,1)\),
\(e_{k+1}=\sqrt L e_k\), and \(0<\mu\le1/L\); the scaled reference
mass is \(L\mu\). Both charges are at most the common ceiling $e_{\mathrm{prim}}$ of the
photon and fermion estimates specified in the theorem below. Require
\(\ell/(LM)\) and \(\ell/(L^2M)\) to be odd integers at least \(13\).
Fix an integer \(r\ge2\). For \(v=-\log e\), put
\[
 r_k=\min\{L^i:i\ge0\text{ an integer},\ L^i\ge v^r\},
 \qquad \ell\ge Mr_k .
\]



An activity in $K_k^{x,\C}$ is a balanced family of Grassmann
polynomials in ordinary and ordered-pair generators, holomorphic on
$\widetilde R_k\times\mathfrak X_k$, local in its polymer and its
connection bonds, and covariant under the lattice, independent local
real-gauge and charge-conjugation actions. The real space $K_k^x$
also satisfies the antilinear reality identity. We use exactly the
domains, pair support and actions of \cite[\S2]{NormalizedContraction},
recalled in \cref{sec:activities}. In particular, $\widetilde R_k$
permits independent local real gauges; it imposes no bound on a real
harmonic part. With $d_M(X)$ the connected-net length divided by $M$,
the weights and norm are
\[
 h_k=(e_k^{-1/4},e_k^{-1/4+\varepsilon}),\qquad
 \|\mathscr E\|_k^x=\sup_X\exp(\kappa d_M(X))
       \sup_{\widetilde R_k\times\mathfrak X_k}\|\mathscr E(X)\|_{h_k}.
\]
The coefficient norm is the sum of absolute ordered coefficients
times their generator weights. Set
\begin{equation}\label{eq:rg-windows}
 w_k^{\mathrm m}=e_k^{3/4-8\varepsilon},\quad
 w_k^{\mathrm E}=e_k^{1/4-7\varepsilon},\quad
 \operatorname{Ball}_k^x(c)=
 \{(m,\mathscr E):|m|\le cw_k^{\mathrm m},\
                 \|\mathscr E\|_k^x\le cw_k^{\mathrm E}\}.
\end{equation}
At $k=0$ the activity is required to be zero. The interaction mass
$m$ is distinct from the positive reference mass $\mu$. Put
$n_k=\ell^3$, $n_{k+1}=(\ell/L)^3$, and $q=e_k$ at $k=0$.
The real photon domain is
$\mathrm{Dom}_j=\{A:H_j^{\mathrm{ph}}A\in\frac14\widetilde R_j\}$.
The density $\mathscr F_j^x$ and normalized block operator $\mathbf T_k^x$
are defined by \eqref{eq:rg-step-input-density} and
\eqref{eq:rg-step-block-transformation}; the latter integrates over
$\mathrm{Box}_k=\{\sup_\Upsilon|\zeta(\Upsilon)|\le p_{0,k}\}$.

The quantitative constants are
\[
 \begin{gathered}
 \bar\varepsilon=1/200,\qquad
 0<\varepsilon\le\bar\varepsilon,\qquad 1-\varepsilon<\alpha<1,\\
 c_X=54,\quad \kappa_0=54\log72,\quad K_{\mathrm{pol}}=\exp(\kappa_0)/36,\quad
 c_{\rm rb}=c_XK_{\mathrm{pol}},\quad c_{\rm fl}=61c_{\rm rb},\\
 c_C=3L^{9/2},\quad b=1,\quad b_L=b/L,\quad b_-=1-L^{-1},\\
 C_T=2\exp(3/2)
       [1+b_-^{-1}(54L\exp(6)+12\exp(1)+1)],\\
 c_S=10^5c_X^5K_{\mathrm{pol}}^5\exp(6\kappa),\qquad
 \kappa_w=\kappa+\kappa_0+3 .
 \end{gathered}
\]
Here $K_{\mathrm{pol}}$ is the scalar polymer-sum constant of
\cite[Appendix~B, (81)]{NormalizedContraction}.
The primitive constants \(C_w,\gamma_w,\alpha_0>0\), the photon
local-gauge constant \(C_B\), and the two exponential-sum bounds
\(c_\gamma\) are those of the photon and fermion families in hypothesis~(i) below; \(C_w,c_\gamma\ge1\).


\begin{theorem}[The exact nonlinear step]\label{thm:rgnl-current-step}
In the setting above, assume the following simultaneous hypotheses.
\begin{enumerate}[label=\textup{(\roman*)},leftmargin=*]
\item Use a single system of photon and fermion primitives on this
torus. Assume the full hypotheses and conclusions of
\cite[Theorems~A--B and Proposition~5.1]{MaxwellBlockManuscript}
and
\cite[Theorem~A, Propositions~4.1,
4.2,
4.3,
6.1 and
6.2]{FermionBlockManuscript}.
Use common constants $(\alpha_0,\gamma_w,\kappa_w;C_w)$ for their
analytic kernel and short H\"older bounds on
$|s_a|\le M^{\alpha_0}$. This assumption includes component
restriction, row vanishing, local read bonds, the exact operators at
$s=1$, and their lattice, gauge, transpose and reality identities.
At $k=0$ retain the stated diagonal families and chamber covariance.
Use the weighted integration conventions of
\cite[Theorem~B and Proposition~5.1]{FermionBlockManuscript}
and the extraction, reconstruction and geometry of
\cite[Theorems~A--B and Appendix~B]{NormalizedContraction}.
The fine inverses required for reconstruction exist at both signs
of every total connection in $\widetilde G_k$ by the fermion hypothesis.

\item The extraction and scaling parameters satisfy
\begin{equation}\label{eq:nlm-normalization-hypotheses}
 \begin{gathered}
 L^{1/8-3\bar\varepsilon/2}\ge2400,\qquad e^\varepsilon\le\tfrac14,\qquad
 160L^{1+3\varepsilon/2}Me^\varepsilon\le1,\\
 9ML^\varepsilon e^{2\varepsilon}\le1,\qquad
 90c_XLMe^\varepsilon\le1,\\
 \kappa\ge7\kappa_0+7,\qquad
 L(\kappa-6\kappa_0-4)-1\ge\kappa .
 \end{gathered}
\end{equation}
\item The weakening radius and constants satisfy
\[
 C_w\ge4\exp(2\gamma_w)M^{27\alpha_0},\qquad
 M^{\alpha_0}-1\ge\exp\big((2L+1)^3\kappa+\log(6c_X)\big).
\]
For the localizations, connected expansion, scaling and determinant
reblocking, assume
\begin{align}
 &p_{0,k}\le e^{-\varepsilon/4},\qquad
       12\exp(-p_{0,k}^2/2)\le e^7,\qquad
       2e^{5\varepsilon/2}\le1,\qquad
       c_{\mathrm{fl}}e^\varepsilon\le1,\qquad3Me^\varepsilon\le1,
                                                        \notag\\
 &2c_CC_w^2c_\gamma^2e^{\varepsilon/4}\le1,\qquad
       16c_Xc_\gamma C_wC_T^2M^3e^{\varepsilon/2}\le1,\qquad
       96M^3e^{1/4}\le1,                                \notag\\
 &32C_wC_Te^{3\varepsilon/2}\le1,\qquad
       88M^3e^{1/4-6\varepsilon}\le1,\qquad
       112C_w^2c_\gamma e^{2\varepsilon}\le3,             \notag\\
 &c_Se^{1/4-5\varepsilon}\le1,\qquad
       c_SL^3e^\varepsilon\le1,\qquad e^\varepsilon\le\tfrac14,\qquad
       8c_{\mathrm{rb}}L^3C_dM^3e^{2\varepsilon}\le1.
                                      \label{eq:rgnl-fixed-smallness}
\end{align}
Here $C_d$ is the local determinant constant in hypothesis~(i).
\item For fixed positive exponents \(p_0,p\), put
\(p_{0,k}=v^{p_0}\), \(p_k=v^p\). The photon box and domain bounds
and the determinant gain require
\[
 \begin{gathered}
 p_{0,k}\le e^{-\varepsilon/4},\qquad
 12\exp(-p_{0,k}^2/2)\le e^7,\qquad
 8C_Bp_k\le e^{-3/4+3\varepsilon},\\
 8e^{3/4-7\varepsilon/2}\le1,\qquad L^{3/4}\ge8/3,\qquad
 v\ge9,\qquad 8M(v/4+4)e^\varepsilon\le1,\\
 e^{1/4}+e^\varepsilon<1,\qquad
 4e^{1/4+2\varepsilon}+2e^{3\varepsilon}\le1 .
 \end{gathered}
\]

\end{enumerate}

There is a local connected logarithm $E_k^\#(m,\mathscr E)$ on
$\widetilde R_k^-\times\mathfrak X_k$ for every complex datum in
$\operatorname{Ball}_k^x(4)$. It has the lattice symmetry with
$LM$ translations, independent local real-gauge covariance and
charge conjugation. Each monomial is balanced; real data give the
reality identity. Its scaled sum
\[
 E_k^*(m,\mathscr E)(Y;A,C)
   =\sum_{\bar X^L=LY}S_LE_k^\#(m,\mathscr E)(X;A_L,C_L)
\]
belongs to $K_{k+1}^{x,\C}$ and has norm at most
$e_k^{1/4-6\varepsilon}$. Here $A_L(b)=L^{-1/2}A(b/L)$;
$S_L$ scales ordinary legs by $L^{-1}$ and pair legs by
$L^{-1-\alpha}$, and $\bar X^L$ is the union of $LM$-cubes
containing $X$. The correction is norm continuous on each affine
complex line's closed allowed set and norm analytic in its interior.
For two data in $\operatorname{Ball}_k^x(1)$,
\[
 \|E_k^*(m_1,\mathscr E_1)-E_k^*(m_2,\mathscr E_2)\|_{k+1}^x
 \le e_k^{1/4-6\varepsilon}
 \left(\frac{|m_1-m_2|}{w_k^{\mathrm m}}+
       \frac{\|\mathscr E_1-\mathscr E_2\|_k^x}{w_k^{\mathrm E}}\right).
\]
The second quotient is omitted at $k=0$.

For the local determinant terms of
Proposition~\extref{6.1}{FermionBlockManuscript}, put
\[
 \delta E_k^{\mathrm{det}}(X;\mathscr A,\mathscr C)
       =E_k^d(X,\mathscr A+\mathscr C)-E_k^d(X,\mathscr C),
 \qquad
 E_k^{\mathrm{det}}(X,\mathscr C)
       =E_k^d(X,\mathscr C)-E_k^d(X,0).
\]
These are scalar activities. We use reblocking
$\mathcal L$ from \exteq{34}{NormalizedContraction}.
The maps $R$, $(\cdot)^{\mathrm{nat}}$, $m(\cdot)$ and
$\varepsilon_{\mathrm v}(\cdot)$ are the subtraction, adjustment and
extraction maps of \cite[Theorem~A]{NormalizedContraction}.
The outgoing maps are
\begin{align}
 \mathscr E^+&=\mathcal L(R\mathscr E)^{\mathrm{nat}}
                        +\mathcal L\delta E_k^{\mathrm{det}}
                        +E_k^*(m,\mathscr E), \label{eq:rgnl-full-step}\\
 m^+&=L(m+m(\mathscr E)),\\
 \varepsilon^{\mathrm v,+}
       &=L^3(\varepsilon^{\mathrm v}
                  +\varepsilon_{\mathrm v}(\mathscr E)
                                      +\delta\varepsilon_k^0),\\
 \mathfrak P_{k+1}(Y,\mathscr C')
       &=\sum_{\bar X^L=LY}
           \big(\mathfrak P_k(X,\mathscr C'_L)
                          +E_k^{\mathrm{det}}(X,\mathscr C'_L)\big).
                                      \label{eq:rgnl-external-scalar-map}
\end{align}
Here $\mathfrak P_k$ is any given family of scalar
functions of the external slot.


Let $\varepsilon^{\mathrm v},m$ be real,
$\mathscr E\in K_k^x$, and
$(m,\mathscr E)\in\operatorname{Ball}_k^x(1)$, with
$\mathscr E=0$ at $k=0$. For $A'\in\mathrm{Dom}_{k+1}$
with arbitrary harmonic part and
$\mathscr C'\in\mathfrak X_{k+1}$, the exact density satisfies
\begin{equation}\label{eq:rgnl-exact-step}
 \mathbf T_k^x\mathscr F_k^x[\varepsilon^{\mathrm v},m,\mathscr E]
   =\mathscr F_{k+1}^x
            [\varepsilon^{\mathrm v,+},m^+,\mathscr E^+],
\end{equation}
with external scalar families $\mathfrak P_k$ and
$\mathfrak P_{k+1}$ on the respective sides.
Its integral satisfies
\begin{equation}\label{eq:rgnl-integrated-step}
 y_{k+1}^x(A';\mathscr C')
  =\int_{\mathrm{Box}_k}\dd\mu_I(\zeta)\,
        y_k^x(A_k^{\min}(A')+\mathrm{Sq}_k\zeta;\mathscr C'_L),
\end{equation}
where $y_j^x=\int\mathrm D\Psi\,\mathscr F_j^x$ with the
data and scalar family of that level. The photon integral uses the
unnormalized Gaussian restriction to the box.
Moreover $\mathscr E^+\in K_{k+1}^x$ and
\begin{align}
 \|\mathscr E^+\|_{k+1}^x
 &\le c_{\mathrm{fl}}L^{-1/4+2\varepsilon}\|\mathscr E\|_k^x
           +e_{k+1}^{1/4-2\varepsilon}+e_k^{1/4-6\varepsilon},
                                      \label{eq:rgnl-activity-step-bound}\\
 |m^+-Lm|
 &\le Lc_{\mathrm{fl}}e_k^{1/2}\|\mathscr E\|_k^x,
                                      \label{eq:rgnl-mass-step-bound}\\
 |\varepsilon^{\mathrm v,+}-L^3\varepsilon^{\mathrm v}|
 &\le L^3(c_{\mathrm{fl}}\|\mathscr E\|_k^x+e_k^7).
                                      \label{eq:rgnl-vacuum-step-bound}
\end{align}
At level zero this formula integrates all $192n_0$
incoming pairs of generators, with chamber measure $1/48$,
and leaves $4n_1$ outgoing pairs. Its data are
\[
 \mathscr E^+=\mathcal L\delta E_0^{\mathrm{det}}+E_0^*(m,0),
 \quad m^+=Lm,\quad
 \varepsilon^{\mathrm v,+}
                         =L^3(\varepsilon^{\mathrm v}+\delta\varepsilon_0^0).
\]
In particular
$0\le\varepsilon^{\mathrm v,+}-L^3\varepsilon^{\mathrm v}
\le L^3q^7$.
All maps commute with setting the external connection to zero.
The estimates are independent of the side and harmonic part, for
one primitive system satisfying the stated hypotheses. The theorem
asserts one step on the two admitted sides; it makes no assertion
that the outgoing datum lies in the next unit ball.
\end{theorem}



The proof develops the exact density and nonlinear arguments of
\cite[\S\S13--14]{QED3Master}. The quadratic integration in \cref{sec:rg-step,subsec:rg-step-first-vertex} first
separates a new fermion density, a determinant quotient and a finite
fluctuation integral. The main difficulty is to express that last
factor by activities with local coefficients. Two finite
localizations in \cref{sec:rg-nonlinear-step} do this: the first resolves the change of connection
and reconstructs the mean field; the second resolves the fermion
fluctuations. Component restriction makes disconnected interpolation
terms vanish before any infinite expansion is taken.

The resulting activity satisfies
\[
 \|F(Y)\|_{(h_k/4;\mathsf w_W)}
 \le36c_X^2K_{\mathrm{pol}}^2e_k^{1/4-5\varepsilon}
                 \exp[-(\kappa-2\kappa_0-2)d_M(Y)].
\]
The connected-polymer expansion in \cref{subsec:rgnl-cluster-logarithm}
converges absolutely, with a
geometric tail in the number of polymers. Its exponential equals
the original finite integral by absolute regrouping. For charge
conjugation, the unbarred fluctuation contains $\Gamma_k$ whereas
the barred one does not. Equality of their transformed Wick
covariances, followed by inversion of a finite subset sum, gives
the symmetry of each local logarithm. Scaling and Cauchy's formula
give the correction and coupling bounds. \Cref{subsec:rgnl-current-step}
evaluates the two finite localization sums and recovers the density
identity with its determinant and photon-box factors.

The extraction and reconstruction arguments are proved separately in
\cite{NormalizedContraction}; their exact hypotheses are retained
above. The elementary labelled-tree degree formula used in the
connected sum is the Pr\"ufer formula in
\cite[Proposition~4]{PostnikovLin2019}. All later sections use the
hypotheses of \cref{thm:rgnl-current-step}. No limiting trajectory is
part of the present theorem.
\section{Activities, connections and coefficient norms}
\label{sec:activities}

We specify the local activity space used in \cref{thm:rgnl-current-step}.
These are the conventions of \cite[\S2]{NormalizedContraction}.

\subsection{Chamber sites, pairs and connection domains}

Write \(O_{\mathrm{cub}}\) for the signed permutation matrices on
\(\R^3\). A chamber site is \(x=(c,g)\), with photon projection
\(\pi x=c\) and \(g\in O_{\mathrm{cub}}\). Integration over chamber
sites uses \(\dd x=(\eta_k^3/48)\sum_{c,g}\). We take
\(M\)-cubes centered at \(M\Z^3/\ell\Z^3\); their lattice sites
form a partition. Blocks of every odd side used below are centered
in the same way.

For the three internal neighbors let \(s_0\) change the sign of
the third coordinate, let \(s_1\) exchange the second and third
coordinates, and let \(s_2\) exchange the first and second.
Let \(r_{e_1}\) change the sign of the first coordinate. Set
\[
 x^{(i)}=(c,gs_i)\quad(0\le i\le2),\qquad
 x^{(3)}=(c+\eta_kge_1,gr_{e_1}).
\]
Each neighbor map is an involution. An oriented bond field
\(B\) is antisymmetric under reversal and periodic on the photon
lattice. Write \(B_\mu(c)=B(c,c+\eta_ke_\mu)\), and use the
forward differences
\[
 \partial_\nu B_\mu(c)=
       \eta_k^{-1}\bigl(B_\mu(c+\eta_ke_\nu)-B_\mu(c)\bigr),
 \qquad
 (d\lambda)_\mu(c)=
       \eta_k^{-1}\bigl(\lambda(c+\eta_ke_\mu)-\lambda(c)\bigr).
\]
Thus \((dB)_{\nu\mu}=\partial_\nu B_\mu-\partial_\mu B_\nu\).
Distances between photon sites are sup distances, using the
shortest lift when a distance bound is imposed. For \(0<\alpha<1\)
the three norms are
\[
 \|B\|_1=\|B\|_\infty,\qquad
 \|B\|_2=\|\partial B\|_\infty,\qquad
 \|B\|_3=\|\delta_\alpha\partial B\|_\infty
 =\sup_{\substack{\mu,\nu,c,c'\\0<d(c,c')\le1}}
 \frac{|\partial_\nu B_\mu(c')-\partial_\nu B_\mu(c)|}
      {d(c,c')^\alpha}.
\]
Local versions use only the indicated cells in the given region.
Let \(\square^\dagger\) be the concentric enlargement of an
\(M\)-cube to side \((2L+1)M\). Put
\(\theta_i=e_k^{-3/4+i\varepsilon}\).
The connection domains are
\begin{equation}\label{eq:connection-domains}
\begin{split}
 R_k&=\{B:\|B\|_i<\theta_i,\ 1\le i\le3\},\\
 \widetilde R_k&=\{B_0+B_1:B_1\in R_k,\ B_0\text{ real, and on each }
          \square^\dagger,\\
 &\hspace{25mm} B_0=B_\square^{\mathrm b}+d\lambda_\square
   \text{ with }B_\square^{\mathrm b},\lambda_\square\text{ real and }
          \|B_\square^{\mathrm b}\|_i<\theta_i\},\\
 \mathfrak X_k&=\{C:\|C\|_\infty<e_k^{-1+2\varepsilon},\
                       \|\partial C\|_\infty<e_k^{-1+3\varepsilon}\}.
\end{split}
\end{equation}
Only existence of these representatives is required.
For \(0<c\le1\), the notation \(c\widetilde R_k\) multiplies
all three bounds for both summands by \(c\).
For \(\widetilde R_k^-\) the local real bounds remain
\(\theta_i\), while the global summand has bounds
\((1-e_k^\varepsilon)\theta_i\).

For a short ordered pair \(x=(c,g),x'=(c',g')\), let
\(\tau B(c,c')\) be the arithmetic mean of the sums of \(B\)
along the six coordinate-order paths from \(c\) to \(c'\).
The bond values in this sum are unscaled; define
\[
 U_B(x,x')=\exp\{ie_k\eta_k\tau B(c,c')\},\qquad
 (\delta_{\alpha,B}\psi)(x,x')
       =\frac{U_B(x,x')\psi(x')-\psi(x)}
                    {\max(|c-c'|_\infty,\eta_k)^\alpha}.
\]
For pairs within one cluster the transport is one. Barred fields
use \(U_B^{-1}\). These formulas specify evaluation; the pair
generators in the algebra remain independent.

We use Pauli matrices and
\[
 \gamma_\mu=\sigma_3\otimes\sigma_\mu,\quad
 \gamma_5=\sigma_1\otimes1_2,\quad \Sigma_3=\sigma_3\otimes1_2,
 \qquad
 \mathsf C=i\gamma_2,\quad
 \varkappa=\gamma_5\mathsf C=\sigma_2\otimes\sigma_2 .
\]
Here \(\mathsf C\) is real with
\(\mathsf C^T=\mathsf C^{-1}=-\mathsf C\), while
\(\varkappa\) is a real symmetric involution. The spin group below
is the group generated by \(\gamma_1,\gamma_2,\gamma_3,\gamma_5\).
Its diagonal coordinate action \(r(S)\) is determined by
\(S\gamma_\mu S^{-1}=r(S)_{\mu\mu}\gamma_\mu\).
Connections transform by pullback on oriented bonds.

\subsection{The local Grassmann algebra}
\label{subsec:rg-activity-spaces}

Fix a level $k$, mesh $\eta_k=L^{-k}$ and an admissible side
$\ell$, with $\ell/(LM)$ odd and at least $13$.
Let $\mathcal P_k$ be the collection of face-connected polymers
of $M$-cubes on this torus. A polymer is small when
$d_M(X)\le L$. A small polymer has an injective lift, unique up
to a period translation, by the diameter and lifting estimates
of Appendix~\extref{B}{NormalizedContraction}.
Write $\operatorname{Vol}(X)=M^3|X|_M$.

The independent Grassmann generators are $\psi$, $\bar\psi$ on
the chamber sites $x=(c,g)$ and $\chi$, $\bar\chi$ on ordered
pairs. For distinct chamber sites on the covering lattice, put
$d_X(x,x')=\max(|c-c'|_\infty,\eta_k)$ and restrict pairs to
$d_X\le M$. Pairs on the torus are simultaneous period classes;
their shortest lifts are unique since $2M<\ell$. A pair is
admissible for $X$ when its photon sites and all six rectilinear
paths between them lie in $X$, or when it is $(x,x^{(3)})$ and
one of its photon sites belongs to $X$.
Thus a pair may leave $X$ only along a bond of type $3$.
Pairs inside a cluster have $d_X=\eta_k$ and are admissible when
their common photon site belongs to $X$.
Ordinary generators acquire the sign $(-1)^{n_1+n_2+n_3}$
under translation by $\ell n$. Pair generators acquire the same
sign when both endpoints are translated by $\ell n$.
The kernels are antiperiodic in each ordinary argument and in
each simultaneous pair argument with these conventions.

Let $\xi=(x,\beta,\omega)$ index an ordinary leg and
$\zeta=(x,x',\beta,\omega)$ a H\"older leg, where
$\omega=0$ is unbarred and $\omega=1$ is barred.
We use the index measures
\[
 \nu(\xi)=\eta_k^3/48,\qquad
 \nu(\zeta)=(\eta_k^3/48)^2.
\]
Integrals sum the spin and bar indices as well. An activity
has a unique expansion with separately antisymmetric kernels,
\begin{equation}\label{eq:rg-activity-kernels}
 \mathscr E(X;\mathscr A,\mathscr C)
 =\sum_{n,m}\frac1{n!m!}\int\dd\xi\,\dd\zeta\,
      \mathscr E_{nm}(X;\mathscr A,\mathscr C;\xi;\zeta)
                       \psi(\xi_1)\cdots\psi(\xi_n)
                       \chi(\zeta_1)\cdots\chi(\zeta_m).
\end{equation}
Barred generators are included in this notation.
The generator weights, as distinct from the index measures, are
\[
 h_{k,1}=e_k^{-1/4},\qquad h_{k,2}=e_k^{-1/4+\varepsilon}.
\]
For an element $F$ the weighted coefficient norm is
\begin{equation}\label{eq:rg-activity-coefficient-norm}
 \|F\|_{h_k}
       =\sum_{n,m}\frac{h_{k,1}^nh_{k,2}^m}{n!m!}
                                   \int\dd\xi\,\dd\zeta\,|F_{nm}|.
\end{equation}
In particular a generator has norm $h_{k,i}$;
the weight $\nu$ enters only when coefficients are written as
integral kernels.

An activity family in $K_k^{x,\C}$ is balanced: every monomial
has as many barred as unbarred legs. Its kernels are holomorphic
on $\widetilde R_k\times\mathfrak X_k$, vanish unless all ordinary
legs lie over $X$ and all pairs are admissible for $X$, and read
the two connections only on bonds meeting $X$. The norm is
\begin{equation}\label{eq:rg-activity-norm}
 \|\mathscr E\|_k^x
  =\sup_{X\in\mathcal P_k}\exp(\kappa d_M(X))
    \sup_{(\mathscr A,\mathscr C)\in\widetilde R_k\times\mathfrak X_k}
                          \|\mathscr E(X;\mathscr A,\mathscr C)\|_{h_k}.
\end{equation}
We require $\kappa\ge\kappa_0$ and finite norm.

The required symmetries act on both the connections and the
generators. For $\tilde g=(S,a)$ with the coordinate sign action
$r(S)$ and $a\in M\Z^3$, put
\[
 \tilde g(c,g)=(r(S)c+a,r(S)g),\qquad
 a_{\tilde g}\psi(x)=S\psi(\tilde g^{-1}x),\qquad
 a_{\tilde g}\bar\psi(x)=(S^{-1})^T\bar\psi(\tilde g^{-1}x).
\]
The same spin matrices act on pair generators after transforming
both endpoints. Local gauge transformations multiply an unbarred
ordinary or pair leg by $\exp(ie_k\Lambda(c))$, using the photon
site of its first point, and a barred leg by the inverse phase.
Write this automorphism as $G_\Lambda$.
Charge conjugation and the antilinear reality map act by
\[
 \mathrm{Cg}\,\psi=\mathsf C\bar\psi,\qquad
 \mathrm{Cg}\,\bar\psi=-\mathsf C\psi,\qquad
 \mathbf K\psi=\varkappa\psi,\qquad
 \mathbf K\bar\psi=\varkappa\bar\psi,
\]
and by the same rules on the H\"older generators.

The covariance requirements on a family are
\begin{equation}\label{eq:rg-activity-symmetries}
 \begin{split}
 a_{\tilde g}\mathscr E(\tilde gX;\tilde g\mathscr A,
                                      \tilde g\mathscr C)
                                  &=\mathscr E(X;\mathscr A,\mathscr C),\\
 \mathrm{Cg}\,\mathscr E(X;-\mathscr A,-\mathscr C)
                                  &=\mathscr E(X;\mathscr A,\mathscr C).
 \end{split}
\end{equation}
For real scalars $\lambda,\lambda'$ and two pairs in the domain
which satisfy $\mathscr A^1=\mathscr A-d\lambda$ and
$\mathscr C^1=\mathscr C-d\lambda'$ on bonds at distance at
most one from $X$, the local gauge requirement is
\[
 G_{\lambda+\lambda'}\mathscr E(X;\mathscr A^1,\mathscr C^1)
                               =\mathscr E(X;\mathscr A,\mathscr C).
\]
The real subspace $K_k^x$ has in addition
\[
 \mathbf K\mathscr E(X;-\overline{\mathscr A},-\overline{\mathscr C})
                               =\mathscr E(X;\mathscr A,\mathscr C).
\]
All requirements refer to the generators before evaluating fields.
They are essential to the form of the counterterms.

The spaces $K_k^\C,K_k$ without the second connection argument
have the same definitions on $\widetilde R_k$. Restriction
$\mathscr E|_0(X;\mathscr A)=\mathscr E(X;\mathscr A,0)$ maps
$K_k^{x,\C}$ into $K_k^\C$ and has norm at most one. The real
subspaces also restrict. Analyticity, locality and the symmetries
follow by setting the second argument and its gauge function to
zero; the norm bound follows from the smaller supremum.

\subsection{Total connections}
\label{subsec:rg-total-connection-domain}

The comparison domain $\widetilde G_k$ consists of fields $B_0+B_1$
with $B_0$ real and $\|B_1\|_\infty<e_k^{-1+\varepsilon}$,
where on each $\square^{\sim5}$ the field $B_0$ has a real-gauge
representative of supremum norm below $e_k^{-1+\varepsilon}$.
Here $\square^{\sim5}$ has the same center and side $11M$.
The real-gauge representatives in the domains defined above may
be chosen independently on each enlarged cube.

Assume $L\ge5$, $0<\varepsilon\le1/200$ and
\begin{equation}\label{eq:rg-domain-charge-conditions}
 e_k^\varepsilon\le\tfrac14,\qquad
 e_k^{1/4}+e_k^\varepsilon\le1,\qquad
 4e_k^{1/4+2\varepsilon}+2e_k^{3\varepsilon}\le1.
\end{equation}
Then
\begin{equation}\label{eq:rg-total-domain-inclusion}
 \widetilde R_k+\mathfrak X_k\subset\widetilde G_k.
\end{equation}
Indeed, write $\mathscr A=A_0+A_1$ with $A_1\in R_k$ and
$A_0$ locally real-gauge equivalent on $\square^\dagger$ to
$A_\square^b$ with the three threshold bounds. Restrict the
representative to $\square^{\sim5}\subset\square^\dagger$.
Its supremum norm is below
$e_k^{-3/4+\varepsilon}\le e_k^{-1+\varepsilon}$, and
\[
 \|A_1+\mathscr C\|_\infty
 <e_k^{-3/4+\varepsilon}+e_k^{-1+2\varepsilon}
 =(e_k^{1/4}+e_k^\varepsilon)e_k^{-1+\varepsilon}
 \le e_k^{-1+\varepsilon}.
\]
This is the required decomposition of the total connection.
Real gauge transformations can absorb any real constant part,
locally on the nonwrapping cubes.

Put $\mathfrak i(B)=e_k\|\operatorname{Im}B\|_\infty$ and
$q_d(B)=e_k\|dB\|_\infty$. The same decomposition gives
\begin{equation}\label{eq:rg-total-imaginary-curvature-bounds}
 \mathfrak i(\mathscr A+\mathscr C)
       \le e_k^{1/4+\varepsilon}+e_k^{2\varepsilon}\le1,
 \qquad
 q_d(\mathscr A+\mathscr C)
       \le4e_k^{1/4+2\varepsilon}+2e_k^{3\varepsilon}\le1.
\end{equation}
For the curvature, every plaquette lies in an enlarged cube,
and $dA_0=dA_\square^b$ there. This costs at most
$2\|\partial A_\square^b\|_\infty$; the global small field
costs another $2\|\partial A_1\|_\infty$, and the external
field costs $2\|\partial\mathscr C\|_\infty$.
The imaginary bound uses only that $A_0$ is real.

The domains are open: the strict inequalities define open
sets for the global summand, which can absorb a sufficiently
small complex perturbation while the local real representative
is kept fixed. The norm-defined sets are invariant under sign
and conjugation; the local real-gauge definitions are also
invariant under real gauges and real constant additions.
Coordinate sign changes and $M$ translations preserve them.
Their derivative and H\"older norms are preserved, with a
possible reversal of a bond and translation of its base point.
Thus \eqref{eq:rg-total-domain-inclusion} applies also at the
signs, conjugates and real gauges used below.



\subsection{The operator and photon data}

The injection and cluster mean are
$(Jv)(c,g)=v(c)$ and $(J^*f)(c)=48^{-1}\sum_gf(c,g)$.
The normalized average $\mathcal Q_k=Q_kJ^*$, its reverse
$\mathcal Q_k^\sharp=\mathcal Q_k(-B)^T$, the one-step average
$Q$, and the chamber operator $D_B$ are given in
\cite[\S1]{FermionBlockManuscript}. Put $b_k=(1-L^{-1})/(1-L^{-k})$ for $k\ge1$. At these levels set
\[
 S_k(B)=(D_B+\mu+b_k\mathcal Q_k^\sharp\mathcal Q_k)^{-1},\quad
 H_k(B)=b_kS_k(B)\mathcal Q_k^\sharp(B),\quad
 D_k(B)=b_k-b_k^2\mathcal Q_k(B)S_k(B)\mathcal Q_k^\sharp(B).
\]
The one-step projection is $P_L(B)=Q^T(-B)Q(B)$. The barred map is $\bar H_k(B)=\mathsf C^{-1}H_k(-B)\mathsf C$.
At $k=0$, $H_0=\bar H_0=1$ and $D_0(B)=D_B+\mu$.
The symbol $\Sspace(a,\ell)$ denotes the four-component spinor
space on the $a$-spaced torus, with measure $a^3$; the superscript
$\mathrm{ch}$ adds chamber labels and divides the measure by $48$.
All spaces have the antiperiodic convention stated above.
For later use, the hopping notation of \cite[\S2.1]{FermionBlockManuscript} is
\[
 R_i(B)f(x)=U_B(x,x^{(i)})f(x^{(i)}),\qquad
 \mathcal L_B=\sum_{i=0}^3w_i(1-R_i(B)),\qquad
 M_i(c,g)=\tfrac{w_i}{2}(\rho+\gamma\cdot g\mathbf n_i),
\]
where $(w_0,w_1,w_2,w_3)=(6,6\sqrt2,6\sqrt2,6)$ and
$(\mathbf n_0,\mathbf n_1,\mathbf n_2,\mathbf n_3)=(e_3,(e_2-e_3)/\sqrt2,(e_1-e_2)/\sqrt2,-e_1)$.
Thus $\widehat d_i(c,g)=-g\mathbf n_i$. On a nearest-neighbor photon bond
write $u_B(c,c')=\exp(ie_k\eta_kB(c,c'))$; this is the transport
$U_B(x,x^{(3)})$ when $c,c'$ are its endpoints.

For photon fields use the Landau reconstruction $H_k^{\mathrm{ph}}$,
the constrained minimum $A_k^{\min}$ and the real coordinate map
$C:V_k\to\mathcal Z$ of
\cite[Theorem~A]{MaxwellBlockManuscript}. Write
$C_k^{1/2}=\widetilde\Delta_k^{-1/2}$ and
$\mathrm{Sq}_k=CC_k^{1/2}$. The coordinate basis $\{\Upsilon\}$ is
orthonormal and has the specified block or face-bond supports.
The law $\mu_I$ is the product standard Gaussian in these coordinates.
Its restriction to $\mathrm{Box}_k$ has probability
$\exp(-\ell^3\delta\varepsilon_k^0)$, with
$0\le\delta\varepsilon_k^0\le12\exp(-p_{0,k}^2/2)$.
The symbols $C$ and $\mathscr C$ always denote, respectively,
this real coordinate map and the external connection.

The common constant $c_\gamma$ bounds both
\[
 \sup_\Upsilon\sum_{\Upsilon'}\exp[-\gamma_wd(\Upsilon,\Upsilon')],
 \qquad \sup_b\sum_{b'}\exp[-\gamma_wd(b,b')],
\]
where the second sum is over unit bonds. The constants $C_w$ and
$\gamma_w$ dominate the photon bounds and the corresponding
fermion bounds as stipulated in hypothesis~(i).
\section{The exact block density}\label{sec:rg-step}

The block transformation has three factors: a new quadratic
fermion density, a determinant quotient, and a finite integral
over photon and fermion fluctuations. We derive this identity
before expanding its last factor. The integration at level
zero includes every chamber component and uses the chamber
Berezin measure. The algebraic identities in this section follow
\cite[\S13.1--13.3]{QED3Master}.

\subsection{Integration spaces and real backgrounds}
\label{subsec:rg-step-real-background}

Fix $k\in\Z_{\ge0}$, $\eta_k=L^{-k}$, odd $L\ge5$, and an
integer power $M$ of $L$ with $M\ge L^2$, with $\ell/(LM)$ odd and at least $13$.
The regulator parameter satisfies $0<\rho<1$.
Keep $b=1$ as a letter, and write
$b_L=b/L$ and
$b_k=b(1-L^{-1})/(1-L^{-k})$ for $k\ge1$.
Write $\mathscr V_k=\Sspace(1,\ell)$ for $k\ge1$, with
four spin components and counting measure, and
$\mathscr V_0=\Sspace^{\mathrm{ch}}(1,\ell)$, with
$192$ components per photon site and measure $1/48$.
The one-step averaging operator into $\Sspace(L,\ell)$ is
\[
 \mathcal A_k(B)=
 \begin{cases}
 Q(B),&k\ge1,\\
 Q(B)J^*,&k=0,
 \end{cases}
 \qquad
 \mathcal A_k^\sharp(B)=\mathcal A_k(-B)^T.
\]
The target measure is $L^3$. On these spaces set
\begin{equation}\label{eq:rg-step-quadratic-operators}
 \mathcal D_k(B)=
 \begin{cases}D_k(B),&k\ge1,\\D_B+\mu,&k=0,\end{cases}
 \qquad
 K_k(B)=\mathcal D_k(B)+b_L\mathcal A_k^\sharp(B)\mathcal A_k(B).
\end{equation}
Thus $K_k=D_k+b_LP_L$ for $k\ge1$ and
$K_0=F_0=D_B+\mu+b_LJP_LJ^*$.
We write $D_0(B)=D_B+\mu$, so that
$\mathcal D_k(B)=D_k(B)$ at every level.
When it exists, $\Gamma_k=K_k^{-1}$, and
$\delta Z_k(B)=\det K_k(B)$.
The transpose refers to the displayed bilinear measures.
For a real $B$, the phases have modulus one and
$\mathcal A_k^\sharp=\mathcal A_k^*$.

The normalized averages give
$Q(B)Q^T(-B)=1$. On a coarse block $B(z')$ the product
is
$L^{-3}\sum_{z\in B(z')}U_B(z',z)U_B(z',z)^{-1}=1$.
At level zero the additional cluster mean contributes
$48/48=1$. Similarly
$\mathcal Q_k\mathcal Q_k^\sharp=1$ for $k\ge1$.
In a unit block the photon count gives
$\eta_k^3\eta_k^{-3}=1$ and the cluster count again
gives $48/48=1$.
For real $B$ the inverse phases in the embeddings are
the complex conjugates of the averaging phases.
The products $\Pi_1=\mathcal Q_k^\sharp\mathcal Q_k$
and $\Pi_L=\mathcal Q_k^\sharp P_L\mathcal Q_k$ are
therefore orthogonal projections; so are $P_L$ and,
at level zero, $JP_LJ^*$. Their scalar spin kernels
commute with every constant spin matrix.

The fine resolvents used below are those of
\cite[\S4.2]{FermionBlockManuscript}. Explicitly, for $k\ge1$ put
\[
 F_\infty(B)=D_B+\mu+b_k\mathcal Q_k^\sharp\mathcal Q_k,
 \quad
 \mathcal B_{k,y}=(b_k+i\gamma_5y)^{-1}(1-P_L)
                  +(b_k+b_L+i\gamma_5y)^{-1}P_L,
\]
and define
\[
 F_y=F_\infty-b_k^2\mathcal Q_k^\sharp\mathcal B_{k,y}\mathcal Q_k,
 \qquad S_{k,y}=F_y^{-1}
\]
where the inverse exists. In particular $S_k=F_\infty^{-1}$.
The reconstruction maps are
$T_k(B)=\mathcal Q_k(B)(b_k^{-1}D_B+1+\mu/b_k)$ and
$\bar T_k(B)=\mathsf C^{-1}T_k(-B)\mathsf C$;
their recovery identities are \cite[Theorem~B(iv)]{NormalizedContraction}.
At zero these reconstruction maps are identities on chamber fields.

\begin{proposition}\label{prop:rg-step-real-inverses}
Let $B$ be real and $\mu>0$. The operators
$D_B+\mu$, $\mathcal D_k(B)$ and $K_k(B)$ are
invertible at every level at which they are defined.
For $k\ge1$,
\begin{equation}\label{eq:rg-step-real-coarse-inverse}
 D_k(B)^{-1}=b_k^{-1}
           +\mathcal Q_k(B)(D_B+\mu)^{-1}\mathcal Q_k^\sharp(B),
 \qquad \operatorname{Re}D_k(B)^{-1}\ge b_k^{-1}.
\end{equation}
Moreover $\operatorname{Re}D_k>0$ and
$\operatorname{Re}K_k>0$ as quadratic forms for
$k\ge1$, whereas $\operatorname{Re}K_0\ge\mu$.
Each $K_k$ satisfies
$K_k^*=\gamma_5K_k\gamma_5$.
For every real $y$, $K_k+i\gamma_5y$ is invertible,
and for $y\ne0$,
\begin{equation}\label{eq:rg-step-real-resolvent}
 \|(K_k+i\gamma_5y)^{-1}\|_{2\to2}\le |y|^{-1}.
\end{equation}
For $k\ge1$ the fine inverse $S_{k,y}$ also
exists for every real $y$.
In particular $\det D_k(0)$ and $\delta Z_k(0)$
are nonzero, including the determinants on
$\mathscr V_0$.
\end{proposition}

\begin{proof}
At a real connection the chamber hops are self-adjoint
unitary involutions. In the proof of
\exteq{6}{FermionBlockManuscript}, reindexing $x$ by $x^{(i)}$
uses
$U_B(x^{(i)},x)=\overline{U_B(x,x^{(i)})}$.
It gives
\[
 (f,\mathcal L_Bf)=\frac12\sum_iw_i
                  \int\dd x\,|R_i(B)f(x)-f(x)|_2^2\ge0.
\]
The direction vector changes sign under the same
involution; hence
$\sum_iw_i(\gamma\cdot\widehat d_i)R_i(B)$
is anti-self-adjoint. The cancellation
$\sum_iw_i\widehat d_i=0$ yields
$D_B^*=\gamma_5D_B\gamma_5$ and
$\operatorname{Re}D_B=(\rho/(2\eta_k))\mathcal L_B\ge0$.
Thus $\operatorname{Re}(D_B+\mu)\ge\mu$.
For $k\ge1$ adding the positive penalty
$b_k\Pi_1$ gives $\operatorname{Re}F_\infty\ge\mu$.
These inequalities imply injectivity and then
invertibility in the finite-dimensional spaces.

Put $G=D_B+\mu$ and $S_k=(G+b_k\mathcal Q_k^\sharp
\mathcal Q_k)^{-1}$. Multiplication of
$b_k^{-1}+\mathcal Q_kG^{-1}\mathcal Q_k^\sharp$
by $D_k=b_k-b_k^2\mathcal Q_kS_k\mathcal Q_k^\sharp$
gives
\[
 1+b_k\mathcal Q_kG^{-1}
      [1-(G+b_k\mathcal Q_k^\sharp\mathcal Q_k)S_k]
                                               \mathcal Q_k^\sharp=1.
\]
This proves \eqref{eq:rg-step-real-coarse-inverse}.
For an invertible operator $O$,
$\operatorname{Re}(u,Ou)=
\operatorname{Re}(Ou,O^{-1}Ou)$.
In particular $\operatorname{Re}O\ge0$ implies
$\operatorname{Re}O^{-1}\ge0$.
As $\mathcal Q_k^\sharp=\mathcal Q_k^*$ here,
the second term of $D_k^{-1}$ has nonnegative real
part. For a nonzero $v$, consequently,
\[
 \operatorname{Re}(v,D_kv)
       =\operatorname{Re}(D_kv,D_k^{-1}D_kv)
                           \ge b_k^{-1}\|D_kv\|_2^2>0.
\]
Adding $b_LP_L$ preserves strict positivity.
At $k=0$ the same conclusion follows directly by
adding $b_LJP_LJ^*$ to $D_B+\mu$.

The projections commute with $\gamma_5$.
It follows first that
$F_\infty^*=\gamma_5F_\infty\gamma_5$ and
$S_k^*=\gamma_5S_k\gamma_5$.
Taking the adjoint in the formula for $D_k$ and using
$\mathcal Q_k^\sharp=\mathcal Q_k^*$ gives
$D_k^*=\gamma_5D_k\gamma_5$. The penalty has the
same symmetry, proving it for $K_k$.
For $k=0$ this is already the symmetry of its three
terms. Hence
\[
 (K_k+i\gamma_5y)^*(K_k+i\gamma_5y)
                                      =K_k^*K_k+y^2,
 \qquad
 \|(K_k+i\gamma_5y)f\|_2^2
                         =\|K_kf\|_2^2+y^2\|f\|_2^2.
\]
For $y\ne0$ this proves injectivity and the inverse
bound; $y=0$ is covered by strict positivity.
For the fine inverse at $k\ge1$, use the block
matrix with diagonal blocks
$b_k+b_LP_L+i\gamma_5y$ and $F_\infty$,
and off-diagonal blocks
$-b_k\mathcal Q_k$ and $-b_k\mathcal Q_k^\sharp$.
The first diagonal block is invertible by its
idempotent decomposition
\exteq{53}{FermionBlockManuscript}.
Its two Schur complements are $F_y$ and
$K_k+i\gamma_5y$. The latter is invertible,
so the block matrix and then $F_y$ are invertible.
This proves existence of $S_{k,y}$.
The asserted normalizing determinants are therefore
nonzero at the real zero connection.
\end{proof}

\subsection{The next quadratic density}
\label{subsec:rg-step-scale-identification}

Assume now that $\ell/L$ is also admissible, and let
$B=B'_L$. At the outgoing level the mass is $L\mu$ and
the charge is $e_{k+1}=L^{1/2}e_k$.
The following identities hold whenever $S_k(B)$ and
$\Gamma_k(B)$ exist for $k\ge1$, or $\Gamma_0(B)$
exists for $k=0$. This invertibility domain is open
and invariant under $B\mapsto-B$: all elementary
entries are entire functions of the bond values,
and the transpose identities are
\begin{equation}\label{eq:rg-step-transpose}
 \mathcal D_k(B)^T=\mathsf C^{-1}\mathcal D_k(-B)\mathsf C,\qquad
 K_k(B)^T=\mathsf C^{-1}K_k(-B)\mathsf C,\qquad
 \Gamma_k(B)^T=\mathsf C^{-1}\Gamma_k(-B)\mathsf C.
\end{equation}
For the first identity reindex each hopping term by
$x^{(i)}$, using
$M_i(x^{(i)})^T=\mathsf C^{-1}M_i(x)\mathsf C$
and the reverse phase at $-B$. Transposition of the
averages reverses their connection in the product
$P_L$, and the scalar kernels commute with
$\mathsf C$. The remaining identities follow by
products and inversion. Invertible matrices have
analytic inverses on their open determinant domains.

Define the fine-to-$L$ average and its embedding by
\[
 \mathcal Q_k^+(B)=
 \begin{cases}Q(B)\mathcal Q_k(B),&k\ge1,\\Q(B)J^*,&k=0,\end{cases}
 \qquad
 \mathcal Q_k^{+,\sharp}(B)=\mathcal Q_k^+(-B)^T.
\]
They obey
$\mathcal Q_k^+\mathcal Q_k^{+,\sharp}=1$.
Set $\bar b_k=b_{k+1}/L$ and
\begin{equation}\label{eq:rg-step-fine-plus}
 F_k^+(B)=D_B+\mu+\bar b_k\mathcal Q_k^{+,\sharp}\mathcal Q_k^+,
 \qquad \Sigma_k^+(B)=F_k^+(B)^{-1}.
\end{equation}
For $k=0$ this is $F_0$ and its inverse $\Gamma_0$.
For $k\ge1$ it is the fine resolvent operator at
$y=0$, since
$\bar b_k=b_kb_L/(b_k+b_L)$.

To verify existence and the identities needed below,
write $\mathcal B_k=(b_k+b_LP_L)^{-1}$ for $k\ge1$
and consider
\[
 \begin{pmatrix}
 b_k+b_LP_L&-b_k\mathcal Q_k\\
 -b_k\mathcal Q_k^\sharp&F_\infty
 \end{pmatrix}.
\]
Its lower-right Schur complement is $K_k$.
Since $P_L$ is idempotent and scalar in spin,
$\mathcal B_k=b_k^{-1}(1-P_L)+(b_k+b_L)^{-1}P_L$.
The other Schur complement is
\[
 F_\infty-b_k^2\mathcal Q_k^\sharp\mathcal B_k\mathcal Q_k
    =D_B+\mu+\frac{b_kb_L}{b_k+b_L}
                                  \mathcal Q_k^\sharp P_L\mathcal Q_k
    =F_k^+.
\]
The two assumed inverses make the block matrix
invertible, and hence $F_k^+$ is invertible.
Equating the upper-left and lower-left blocks of
its inverse gives
\begin{equation}\label{eq:rg-step-block-inverse-identities}
 \Gamma_k=\mathcal B_k+
       b_k^2\mathcal B_k\mathcal Q_k\Sigma_k^+
                                     \mathcal Q_k^\sharp\mathcal B_k,
 \qquad
 H_k\Gamma_k=b_k\Sigma_k^+\mathcal Q_k^\sharp\mathcal B_k.
\end{equation}
This is the same Schur calculation as
Proposition~\extref{4.2}{FermionBlockManuscript}, retaining
the off-diagonal block for the mean field.

Let
\begin{equation}\label{eq:rg-step-unscaled-coarse}
 D_{k+1}^0(B)=b_L-
                       b_L^2\mathcal A_k(B)\Gamma_k(B)\mathcal A_k^\sharp(B).
\end{equation}
For $k\ge1$, $QP_L=Q$ and $P_LQ^T(-B)=Q^T(-B)$
give
$Q\mathcal B_k=(b_k+b_L)^{-1}Q$ and
$\mathcal B_kQ^T(-B)=(b_k+b_L)^{-1}Q^T(-B)$.
Insert these in \eqref{eq:rg-step-block-inverse-identities}
to obtain
\begin{equation}\label{eq:rg-step-compressed-identities}
 \begin{split}
 D_{k+1}^0(B)
       &=\bar b_k-\bar b_k^2
                            \mathcal Q_k^+\Sigma_k^+\mathcal Q_k^{+,\sharp},\\
 b_LH_k(B)\Gamma_k(B)\mathcal A_k^\sharp(B)
       &=\bar b_k\Sigma_k^+\mathcal Q_k^{+,\sharp}.
 \end{split}
\end{equation}
For $k=0$ both formulas follow directly from
$\bar b_0=b_L$, $H_0=1$,
$\Sigma_0^+=\Gamma_0$, and
$\mathcal Q_0^+=\mathcal A_0$.

Let $\mathsf R_L$ denote the scaling
$f_L(c,g)=L^{-1}f(c/L,g)$ on chamber spinors and
$v_L(z')=L^{-1}v(z'/L)$ on unit spinors.
These are bijections between the corresponding
incoming and outgoing spaces. Scaling preserves
the label and the nested block centers, and
$e_k\eta_kL^{-1/2}=e_{k+1}\eta_{k+1}$ preserves
each transport phase. Consequently
\[
 (D_B+\mu)\mathsf R_L
       =L^{-1}\mathsf R_L(D_{B'}+L\mu),\qquad
 \mathcal Q_k^+(B)\mathsf R_L
       =\mathsf R_L\mathcal Q_{k+1}(B'),
\]
and the corresponding embedding identity holds.
For the average, its $L^{-3}$ factor cancels the
$L^3$ change of the chamber index measure.
Thus
\begin{equation}\label{eq:rg-step-fine-scaling}
 F_k^+(B)\mathsf R_L
                   =L^{-1}\mathsf R_LF_\infty(B'),\qquad
 \Sigma_k^+(B)\mathsf R_L
                         =L\mathsf R_LS_{k+1}(B').
\end{equation}
In particular the outgoing inverse exists.
Equations \eqref{eq:rg-step-compressed-identities} and
\eqref{eq:rg-step-fine-scaling} now imply
\begin{equation}\label{eq:rg-step-coarse-scaling}
 D_{k+1}^0(B)\mathsf R_L
                       =L^{-1}\mathsf R_LD_{k+1}(B'),\qquad
 b_LH_k(B)\Gamma_k(B)\mathcal A_k^\sharp(B)\mathsf R_L
                       =\mathsf R_LH_{k+1}(B').
\end{equation}
All expressions in the second identity act on the
outgoing four-component field.

For independent outgoing Grassmann fields define
\begin{equation}\label{eq:rg-step-mean-fields}
 \Psi^c=b_L\Gamma_k(B)\mathcal A_k^\sharp(B)\Psi'_L,\qquad
 \bar\Psi^c=b_L\Gamma_k(B)^T\mathcal A_k(B)^T\bar\Psi'_L.
\end{equation}
Then, with $\psi'=H_{k+1}(B')\Psi'$ and
$\bar\psi'=\bar H_{k+1}(B')\bar\Psi'$,
\begin{equation}\label{eq:rg-step-induced-fields}
 H_k(B)\Psi^c=\psi'_L,\qquad
 \bar H_k(B)\bar\Psi^c=\bar\psi'_L.
\end{equation}
The first relation is \eqref{eq:rg-step-coarse-scaling}.
For the second use \eqref{eq:rg-step-transpose}
and apply the first relation at $-B$ to
$\mathsf C\bar\Psi'$; the constant matrix
$\mathsf C$ commutes with scaling and the averages.
The covariant H\"older quotients of these induced
fields are also the scaled outgoing quotients:
the transports agree and $d_X$ is divided by $L$,
so the quotient carries the factor $L^{-1-\alpha}$.
This holds for every pair, including a pair inside
one cluster.
Applying $T_kH_k=1$ and its barred version shows
that the means in \eqref{eq:rg-step-mean-fields}
are also $T_k(B)\psi'_L$ and $\bar T_k(B)\bar\psi'_L$.
At level zero these maps are identities on
the chamber field.

The quadratic term in the Gaussian block formula is
therefore exactly
\[
 \langle\bar\Psi'_L,D_{k+1}^0(B)\Psi'_L\rangle_L
                         =\langle\bar\Psi',D_{k+1}(B')\Psi'\rangle_1,
\]
since the pairings of scaled fields multiply by $L$.
Together with \exteq{73}{FermionBlockManuscript},
this identifies the entire outgoing quadratic density.

\begin{proposition}\label{prop:rg-step-determinant-normalization}
On the invertibility domain just specified,
\begin{equation}\label{eq:rg-step-determinant-normalization}
 \det D_{k+1}(B')=
           b^{4n_{k+1}}\frac{\det\mathcal D_k(B)}{\delta Z_k(B)},
 \qquad n_k=\ell^3,\quad n_{k+1}=(\ell/L)^3.
\end{equation}
In particular, with $c_k=b^{-4n_{k+1}}$,
\begin{equation}\label{eq:rg-step-zero-normalization}
 \frac{c_k}{\det D_k(0)}
              =\frac{1}{\det D_{k+1}(0)\delta Z_k(0)}.
\end{equation}
\end{proposition}

\begin{proof}
Apply Sylvester's determinant identity to
$b_L\mathcal A_k\Gamma_k$ and $\mathcal A_k^\sharp$
on the two different spaces. It gives
\[
 b_L^{-4n_{k+1}}\det D_{k+1}^0(B)
       =\det(1-b_L\mathcal A_k^\sharp\mathcal A_k\Gamma_k)
       =\det(\mathcal D_k\Gamma_k)
       =\frac{\det\mathcal D_k(B)}{\delta Z_k(B)}.
\]
The incoming determinant at $k=0$ has size $192n_0$;
the exponent of $b_L$ is the outgoing dimension
$4n_1$, as in every other step.
Equation \eqref{eq:rg-step-coarse-scaling} gives
$\det D_{k+1}^0(B)=L^{-4n_{k+1}}\det D_{k+1}(B')$.
Using $Lb_L=b$ proves
\eqref{eq:rg-step-determinant-normalization}.
At the zero connection all denominators are nonzero
by \cref{prop:rg-step-real-inverses}, so
\eqref{eq:rg-step-zero-normalization} follows.
\end{proof}

\subsection{The density with an external connection}
\label{subsec:rg-step-density}

For this subsection assume $0<\mu\le1/L$ and
$e_{k+1}=L^{1/2}e_k\le e_{\mathrm{prim}}$, with the hypotheses
of Propositions~\extref{3.3}{FermionBlockManuscript},
\extref{4.2}{FermionBlockManuscript} and Theorem
\extref{B}{MaxwellBlockManuscript}. We also impose
\exteq{43}{MaxwellBlockManuscript} and
\eqref{eq:rg-domain-charge-conditions}.
The inverse domain $\widetilde G_k$ is the domain
$\widetilde{\mathcal G}_k$ of the propagator estimates.
Let $A'\in\mathrm{Dom}_{k+1}$ be a real block field,
with arbitrary harmonic part, and let
$\mathscr C'\in\mathfrak X_{k+1}$.
Put
\[
 \mathscr A'=H_{k+1}^{\mathrm{ph}}A',\qquad
 \mathscr A=\mathscr A'_L,\qquad
 \mathscr C=\mathscr C'_L,\qquad
 \mathscr T'=\mathscr A'+\mathscr C',\qquad
 \mathscr T=\mathscr T'_L.
\]
The photon block objects are
\begin{equation}\label{eq:rg-step-photon-fields}
 \begin{split}
 A_k(\zeta)&=A_k^{\min}(A')+\mathrm{Sq}_k\zeta,\\
 \mathscr Z(\zeta)&=H_k^{\mathrm{ph}}\mathrm{Sq}_k\zeta,\qquad
 \mathscr B_k=H_k^{\mathrm{ph}}A_k^{\min}(A')
                                      =\mathscr A-d\omega_k.
 \end{split}
\end{equation}
Here $\omega_k$ is real and periodic, as in
\exteq{30}{MaxwellBlockManuscript}.
The truncated photon measure has the normalization
\begin{equation}\label{eq:rg-step-box-normalization}
 \mu_I(\mathrm{Box}_k)=\exp(-n_k\delta\varepsilon_k^0),\qquad
 \widehat\mu_k=
              \frac{\mathbf1_{\mathrm{Box}_k}}{\mu_I(\mathrm{Box}_k)}\mu_I,
\end{equation}
proved in \exteq{24}{MaxwellBlockManuscript}.
Its denominator is positive: $e_k<1$ makes
$p_{0,k}=(-\log e_k)^{p_0}>0$, and the box has
nonempty interior in the photon coordinate space.

All connections read by the step have their required
domains. Indeed, scaling gives
$\mathscr A\in\frac14L^{-3/4}\widetilde R_k
\subset\widetilde R_k^-$ and
$\mathscr C\in\mathfrak X_k$.
Theorem~\extref{B}{MaxwellBlockManuscript}
gives $A_k(\zeta)\in\mathrm{Dom}_k$ for every
$\zeta\in\mathrm{Box}_k$.
The real-gauge identity in
\eqref{eq:rg-step-photon-fields} yields
\[
 H_k^{\mathrm{ph}}A_k(\zeta)
                 =\mathscr A+\mathscr Z(\zeta)-d\omega_k.
\]
Since $\frac14\widetilde R_k$ is invariant under
real gauges, $\mathscr A+\mathscr Z(\zeta)$ also
lies in this domain. Thus both connection pairs
$(\mathscr A,\mathscr C)$ and
$(\mathscr A+\mathscr Z,\mathscr C)$ lie in
$\widetilde R_k\times\mathfrak X_k$.
By \eqref{eq:rg-total-domain-inclusion},
$\mathscr T$, $\mathscr T+\mathscr Z$ and
$\mathscr T+\mathscr Z-d\omega_k$, and their negatives,
belong to $\widetilde G_k$.
The fine and coarse inverse estimates consequently
apply to them. The outgoing operators at
$\mathscr T'$ then exist by
\eqref{eq:rg-step-fine-scaling}.

For a total connection $B$ let $\operatorname{ev}_B$
insert
\[
 \psi=H_k(B)\Psi,\qquad
 \bar\psi=\bar H_k(B)\bar\Psi,\qquad
 \chi=\delta_{\alpha,B}H_k(B)\Psi,\qquad
 \bar\chi=\delta_{\alpha,-B}\bar H_k(B)\bar\Psi.
\]
At $k=0$ the two field maps are identities.
For complex $m,\varepsilon^{\mathrm v}$ and
$\mathscr E\in K_k^{x,\C}$, with $\mathscr E=0$
at $k=0$, define
\begin{equation}\label{eq:rg-step-interaction-action}
 \begin{split}
 I_k^x[m,\mathscr E](A,C;\Psi)
     &=\operatorname{ev}_{A+C}
         \left(-m\langle\bar\psi,\psi\rangle
                               +\sum_X\mathscr E(X;A,C)\right),\\
 \mathrm{Act}_k^x[m,\mathscr E](A,C;\Psi)
     &=-\langle\bar\Psi,D_k(A+C)\Psi\rangle
                                  +I_k^x[m,\mathscr E](A,C;\Psi).
 \end{split}
\end{equation}
The mass pairing is always the chamber pairing
$\eta_k^3\,48^{-1}\sum_{c,g,\beta}\bar\psi_\beta(c,g)\psi_\beta(c,g)$.
The kinetic pairing uses the measure of $\mathscr V_k$.
For an arbitrary finite family of scalar functions
$\mathfrak P_k(X,C)$ on $\mathfrak X_k$, set
\begin{equation}\label{eq:rg-step-input-density}
 \begin{split}
 \mathscr F_k^x[\varepsilon^{\mathrm v},m,\mathscr E](A,C;\Psi)
   =\det D_k(0)^{-1}\exp\left\{
         \mathrm{Act}_k^x[m,\mathscr E](H_k^{\mathrm{ph}}A,C;\Psi)
                   -\varepsilon^{\mathrm v}n_k
                              +\sum_X\mathfrak P_k(X,C)\right\}.
 \end{split}
\end{equation}
The arguments here are $A\in\mathrm{Dom}_k$ and
$C\in\mathfrak X_k$. The fields $\Psi$ and $\bar\Psi$
are independent Grassmann variables.

Define the block transformation on an even density
$F$ whose coefficients are integrable along
$A_k(\zeta)$ by
\begin{equation}\label{eq:rg-step-block-transformation}
 \begin{split}
 (\mathbf T_k^xF)(A',\mathscr C';\Psi')
  &=c_k\int_{\mathrm{Box}_k}\dd\mu_I(\zeta)
                       \int\mathrm D\Psi\,e^{\mathcal Q_{\mathrm{bl}}}
                                         F(A_k(\zeta),\mathscr C;\Psi),\\
 \mathcal Q_{\mathrm{bl}}
  &=-b_L\big\langle
       e^{-ie_k\omega_k}\bar\Psi'_L
                    -\mathcal A_k(-\mathscr B_k-\mathscr C)\bar\Psi,\,
       e^{ie_k\omega_k}\Psi'_L
                    -\mathcal A_k(\mathscr B_k+\mathscr C)\Psi
                                                \big\rangle_L.
 \end{split}
\end{equation}
The gauge function in the two outgoing phases is
evaluated at the $L$-spaced sites.
The Berezin functional is the weighted functional
of \exteq{65}{FermionBlockManuscript}; its incoming mode
measure is $1$ for $k\ge1$ and $1/48$ for $k=0$.

Its normalization is
\begin{equation}\label{eq:rg-step-block-marginal}
 \int\mathrm D\Psi'\,(\mathbf T_k^xF)(A',\mathscr C';\Psi')
       =\int_{\mathrm{Box}_k}\dd\mu_I(\zeta)
                       \int\mathrm D\Psi\,F(A_k(\zeta),\mathscr C;\Psi).
\end{equation}
To prove this, integrate the outgoing family first.
Reciprocal multiplication of its barred and unbarred
generators removes the two phases and leaves the top
coefficient unchanged. Translation by the two averaged
incoming fields likewise leaves it unchanged: every
translated term that contains an external odd field
has lower degree in the outgoing family.
The remaining outgoing Gaussian exponent is
$-b_L\langle\bar\Psi'_L,\Psi'_L\rangle_L
=-b\langle\bar\Psi',\Psi'\rangle_1$.
Its integral is $b^{4n_{k+1}}=c_k^{-1}$ by
\exteq{68}{FermionBlockManuscript}.
The two Berezin integrations commute, since both
extract an even top-degree coefficient; integration
of the finite-dimensional coefficient functions in
$\zeta$ also commutes with them. This proves
\eqref{eq:rg-step-block-marginal}.

For the density \eqref{eq:rg-step-input-density}
the coefficient integrability follows from the domain
statements above. The bond entries of all inverse
operators are analytic on their invertibility domains,
and the activity kernels are holomorphic. Composing
them with the affine map $\zeta\mapsto
H_k^{\mathrm{ph}}A_k(\zeta)$ gives continuous
coefficients on the compact photon box. Their even
Grassmann exponentials are also continuous there.

\begin{proposition}\label{prop:rg-step-density-identity}
Let the data and domains be as above. Put
\[
 m'=m+m(\mathscr E),\qquad \mathscr E^r=R\mathscr E,\qquad
 m^+=Lm',\qquad
 \varepsilon^{\mathrm v,+}
       =L^3\big(\varepsilon^{\mathrm v}
              +\varepsilon_{\mathrm v}(\mathscr E)+\delta\varepsilon_k^0\big).
\]
At $k=0$ the activity and its extracted terms vanish.
Define the outgoing inserted fields
$\psi'=H_{k+1}(\mathscr T')\Psi'$,
$\bar\psi'=\bar H_{k+1}(\mathscr T')\bar\Psi'$ and
their covariant H\"older quotients, and let
$\psi^{0,\#}$ denote their scaled fine fields,
including the barred ones.
Let the means $\Psi^c,\bar\Psi^c$ be those of
\eqref{eq:rg-step-mean-fields} at $B=\mathscr T$,
and put
\[
 \Psi^f=\Psi^c+\Gamma_k(\mathscr T)W,\qquad
 \bar\Psi^f=\bar\Psi^c+\bar W.
\]
With $\mathscr Z=H_k^{\mathrm{ph}}\mathrm{Sq}_k\zeta$,
set
\begin{equation}\label{eq:rg-step-fluctuation-exponent}
 \begin{split}
 V_k^x={}&
       \mathrm{Act}_k^x[m',\mathscr E^r]
                    (\mathscr A+\mathscr Z,\mathscr C;\Psi^f)
       -\mathrm{Act}_k^x[m',\mathscr E^r]
                    (\mathscr A,\mathscr C;\Psi^f)\\
       &+I_k^x[m',\mathscr E^r](\mathscr A,\mathscr C;\Psi^f)
       -I_k^x[m',\mathscr E^r](\mathscr A,\mathscr C;\Psi^c),\\
 \Xi_k^x={}&\int\dd\mu_I(W)\,\dd\widehat\mu_k(\zeta)\,e^{V_k^x}.
 \end{split}
\end{equation}
Then the exact identity is
\begin{equation}\label{eq:rg-step-density-identity}
 \begin{split}
 &(\mathbf T_k^x\mathscr F_k^x
                [\varepsilon^{\mathrm v},m,\mathscr E])
                                      (A',\mathscr C';\Psi')\\
 &\quad=\det D_{k+1}(0)^{-1}
       \exp\Big\{-\langle\bar\Psi',D_{k+1}(\mathscr T')\Psi'\rangle_1
                  -m^+\langle\bar\psi',\psi'\rangle
                  -\varepsilon^{\mathrm v,+}n_{k+1}\\
 &\hspace{48mm}
          +\sum_X\mathscr E^r(X;\mathscr A,\mathscr C;\psi^{0,\#})
          +\sum_X\mathfrak P_k(X,\mathscr C)\Big\}
                     \frac{\delta Z_k(\mathscr T)}{\delta Z_k(0)}\Xi_k^x.
 \end{split}
\end{equation}
All integrals in this formula are finite coefficient
integrals. Its right side depends on $A'$ through
$H_{k+1}^{\mathrm{ph}}A'$.
\end{proposition}

\begin{proof}
Extraction gives
\[
 I_k^x[m,\mathscr E](A,C;\Psi)
      =-\varepsilon_{\mathrm v}(\mathscr E)n_k
                                      +I_k^x[m',\mathscr E^r](A,C;\Psi).
\]
This is the image of
\exteq{18}{NormalizedContraction} under the
field evaluation homomorphism. At $k=0$ it holds
with zero activity and $m'=m$.

For a fixed $\zeta$, change the integrated variables
by $\Psi\mapsto e^{ie_k\omega_k}\Psi$ and
$\bar\Psi\mapsto e^{-ie_k\omega_k}\bar\Psi$.
This leaves the weighted Berezin functional unchanged,
since each mode's factors cancel in its top monomial,
also for the $192n_0$ modes of weight $1/48$.
The transport identity
\[
 U_{B-d\omega_k}(c,c')e^{ie_k\omega_k(c')}
                       =e^{ie_k\omega_k(c)}U_B(c,c')
\]
intertwines the Dirac operator and every average.
Products and inverses preserve this intertwining,
so it also holds for $H_k$ and $D_k$.
Its covariant difference identity uses the phase
at the first photon site, including for a pair
inside one cluster. The barred identity has the
inverse phase. Gauge covariance of the activity
and invariance of the mass pairing consequently
turn its first connection argument
$\mathscr A+\mathscr Z-d\omega_k$ into
$\mathscr A+\mathscr Z$.
The same change removes the gauge from the block
Gaussian factor. The integral in
\eqref{eq:rg-step-block-transformation} is now
\[
 \int\mathrm D\Psi\,
       \exp\left\{-b_L
          \langle\bar\Psi'_L-\mathcal A_k(-\mathscr T)\bar\Psi,\,
                  \Psi'_L-\mathcal A_k(\mathscr T)\Psi\rangle_L
            +\mathrm{Act}_k^x[m',\mathscr E^r]
                    (\mathscr A+\mathscr Z,\mathscr C;\Psi)\right\},
\]
multiplied by
$\det D_k(0)^{-1}
\exp\{-(\varepsilon^{\mathrm v}+
\varepsilon_{\mathrm v}(\mathscr E))n_k+
\sum_X\mathfrak P_k(X,\mathscr C)\}$.

Separate the background quadratic term by writing
\[
 \begin{split}
 \mathcal R(\Psi)={}&I_k^x[m',\mathscr E^r](\mathscr A,\mathscr C;\Psi)\\
  &+\mathrm{Act}_k^x[m',\mathscr E^r]
                    (\mathscr A+\mathscr Z,\mathscr C;\Psi)
   -\mathrm{Act}_k^x[m',\mathscr E^r]
                    (\mathscr A,\mathscr C;\Psi).
 \end{split}
\]
The last exponent above is then the background
Gaussian exponent plus $\mathcal R(\Psi)$.
Translation by the means \eqref{eq:rg-step-mean-fields}
and the block formula \exteq{71}{FermionBlockManuscript}
give
\[
 \delta Z_k(\mathscr T)
   e^{-\langle\bar\Psi',D_{k+1}(\mathscr T')\Psi'\rangle_1}
          \int\dd\mu_I(W)\,e^{\mathcal R(\Psi^f)}.
\]
The change of variables in this formula uses the
same mode measure for $\Psi$ and $W$.
At $k=0$ it is the identity
\[
 \int\mathrm D\Psi\,e^{-\langle\bar\Psi,F_0(\mathscr T)\Psi\rangle^{\mathrm{ch}}}
       G(\bar\Psi,\Psi)
   =\det F_0(\mathscr T)
          \int\dd\mu_I(W)\,
                    G(\bar W,\Gamma_0(\mathscr T)W),
\]
on the entire $192n_0$-mode space.
The Gaussian measure is
$e^{-\langle\bar W,W\rangle^{\mathrm{ch}}}\mathrm DW$,
with the weight $\sqrt{48}$ in its coefficient bound.

By definition of $V_k^x$,
$\mathcal R(\Psi^f)=
I_k^x[m',\mathscr E^r](\mathscr A,\mathscr C;\Psi^c)+V_k^x$.
Its first summand is independent of $W$ and $\zeta$,
so its even exponential factors out.
The induced-field identity
\eqref{eq:rg-step-induced-fields} gives
\[
 I_k^x[m',\mathscr E^r](\mathscr A,\mathscr C;\Psi^c)
      =-m'\langle\bar\psi'_L,\psi'_L\rangle
                 +\sum_X\mathscr E^r(X;\mathscr A,\mathscr C;\psi^{0,\#}).
\]
The chamber measure changes by $L^3$ and each field
by $L^{-1}$. Hence
$\langle\bar\psi'_L,\psi'_L\rangle
=L\langle\bar\psi',\psi'\rangle$, giving $m^+=Lm'$.
Normalizing the photon-box restriction in
\eqref{eq:rg-step-box-normalization} adds
$-\delta\varepsilon_k^0 n_k$ to the vacuum term.
Since $n_k=L^3n_{k+1}$, this term has exactly
the stated outgoing vacuum coefficient.
Finally \eqref{eq:rg-step-zero-normalization}
turns the remaining constant factor into
$\det D_{k+1}(0)^{-1}\delta Z_k(0)^{-1}$,
proving \eqref{eq:rg-step-density-identity}.

The fluctuation exponent is even and continuous
on the photon box, by the same coefficient argument
as for the input density. Thus its integral is defined,
and its two finite integrations commute.
Every background in the final expression is formed
from $\mathscr A'$ and $\mathscr C'$; the remaining
photon maps and measure are independent of $A'$.
The minimum and the real gauge have disappeared
from that expression, proving the last assertion.
\end{proof}

\subsection{Local substitution of the mean field}
\label{subsec:rg-step-local-substitution}

We recall the result of \cite[Theorem~B(iii)--(iv)]{NormalizedContraction}
that is used between the two localizations. For $k\ge1$ and unit
generators over $Y$, the substitution fixes the old ordinary and
pair generators and sets
\begin{equation}\label{eq:rg-step-local-substitution}
 \operatorname{sub}_{Y,B}\Psi=T_Y^{(1)}(B)\psi+T_Y^{(2)}(B)\chi,
 \qquad
 \operatorname{sub}_{Y,B}\bar\Psi=
              \bar T_Y^{(1)}(B)\bar\psi+\bar T_Y^{(2)}(B)\bar\chi.
\end{equation}
The row measures are $\eta_k^3/48$ for ordinary columns and
$(\eta_k^3/48)^2$ for pair columns. The rows add only ordinary
legs over $Y$ and boundary pairs $(x,x^{(3)})$ admissible for $Y$;
they read only bonds meeting $Y$. The substitution preserves
balance, parity and the lattice, local gauge, charge and reality
identities of the old generators. In particular,
\begin{equation}\label{eq:rg-step-substitution-covariances}
 \begin{split}
 a_g\operatorname{sub}_{gY,gB}
                        &=\operatorname{sub}_{Y,B}a_g,\\
 G_\lambda\operatorname{sub}_{Y,B^1}
                        &=\operatorname{sub}_{Y,B}G_\lambda,
                 \qquad B^1=B-d\lambda\text{ on bonds meeting }Y,\\
 \mathrm{Cg}\operatorname{sub}_{Y,-B}
                        &=\operatorname{sub}_{Y,B}\mathrm{Cg},\\
 \mathbf K\operatorname{sub}_{Y,-\overline B}
                        &=\operatorname{sub}_{Y,B}\mathbf K.
 \end{split}
\end{equation}
The local gauge equality need hold only on the bonds meeting $Y$.
For $B\in\widetilde R_k+\mathfrak X_k$, the row bound gives
\begin{equation}\label{eq:rg-step-substitution-bound}
 \|\operatorname{sub}_{Y,B}F\|_{w_1,w_2}
       \le\|F\|_{w_1,w_2;w_3}
       \quad\text{if }C_T(w_1+w_2)\le w_3.
\end{equation}
The bounds on imaginary part and curvature in
\eqref{eq:rg-total-imaginary-curvature-bounds} verify the hypotheses
of that result. The primitive inverse domain at $B$ and $-B$
also verifies its additional recovery hypothesis. Therefore, on
evaluation at $\psi=H_k(B)\Phi$, $\chi=\delta_{\alpha,B}\psi$ and
the barred fields, the substitution sends $\Psi,\bar\Psi$ to
$\Phi,\bar\Phi$. At zero no such substitution is needed.

\section{The first chamber step}\label{subsec:rg-step-first-vertex}

At $k=0$ write the charge as $q$.
There is no input activity, and $H_0(B)=\bar H_0(B)=1$.
Consequently the exponent in
\eqref{eq:rg-step-fluctuation-exponent} is
\begin{equation}\label{eq:rg-step-first-exponent}
 V_0^x=
  -\langle\bar\Psi^f,(D_{\mathscr T+\mathscr Z}
                                    -D_{\mathscr T})\Psi^f\rangle^{\mathrm{ch}}
  -m\left(\langle\bar\Psi^f,\Psi^f\rangle^{\mathrm{ch}}
                   -\langle\bar\Psi^c,\Psi^c\rangle^{\mathrm{ch}}\right).
\end{equation}
The terms with $\mu$ and $m$ cancel in the connection
difference of the actions; the second difference
of interactions gives the displayed mass term.
Here $\Psi^c=\psi'_L$, $\bar\Psi^c=\bar\psi'_L$,
$\Gamma_0(\mathscr T)=F_0(\mathscr T)^{-1}$,
and $\Psi^f=\Psi^c+\Gamma_0(\mathscr T)W$.
All these are fields in $\mathscr V_0=\Sspace^{\mathrm{ch}}(1,\ell)$.
The identity measure in $W$ has mode measure $1/48$,
and its two-point kernel is $48$ times the identity
matrix of the modes.

For a bond fluctuation $Z$ and an $M$-cube $\square$,
define its local vertex by
\begin{equation}\label{eq:rg-step-first-vertex}
 \begin{split}
 \mathcal V_0(Z,\square)=\frac12\int_\square\dd x\,
   \Big\{&
    \bar\psi(x)^TM_3(x)(e^{iqZ(c,c_3)}-1)
                                  [\psi(x)+\chi(x,x^{(3)})]\\
    &+[\bar\psi(x)+\bar\chi(x,x^{(3)})]^T
               M_3(x^{(3)})(e^{-iqZ(c,c_3)}-1)\psi(x)\Big\},
 \end{split}
\end{equation}
where $x=(c,g)$, $c_3=c+ge_1$ and
$\int_\square\dd x=48^{-1}\sum_{c\in\square}\sum_g$.
The pair $(x,x^{(3)})$ has $d_X=1$ and is admissible
for $\square$ of the second kind. Thus the vertex
has ordinary legs over the cube and precisely the
boundary pair legs allowed in the activity algebra.
Its coefficients are entire functions of $Z$ on
the bonds meeting the cube.

\begin{proposition}\label{prop:rg-step-first-vertex}
Evaluate $\chi=\delta_{\alpha,\mathscr T}f$,
$\bar\chi=\delta_{\alpha,-\mathscr T}\bar f$,
$\psi=f$ and $\bar\psi=\bar f$ at any pair of
antiperiodic odd chamber fields. Then
\begin{equation}\label{eq:rg-step-first-vertex-evaluation}
 \operatorname{ev}\sum_\square\mathcal V_0(Z,\square)
          =-\langle\bar f,(D_{\mathscr T+Z}-D_{\mathscr T})
                                                    f\rangle^{\mathrm{ch}}.
\end{equation}
The two halves have equal images after summing
over cubes. The vertex has the lattice, real-gauge,
charge-conjugation and reality symmetries
\begin{equation}\label{eq:rg-step-first-vertex-symmetries}
 \begin{split}
 a_g\mathcal V_0(gZ,g\square)&=\mathcal V_0(Z,\square),\\
 G_\lambda\mathcal V_0(Z,\square)&=\mathcal V_0(Z,\square),\\
 \mathrm{Cg}\mathcal V_0(-Z,\square)&=\mathcal V_0(Z,\square),\\
 \mathbf K\mathcal V_0(-\overline Z,\square)&=\mathcal V_0(Z,\square).
 \end{split}
\end{equation}
For any generator weights $(w_1,w_2)$,
\begin{equation}\label{eq:rg-step-first-vertex-bound}
 \|\mathcal V_0(Z,\square)\|_{w_1,w_2}
       \le12(1+\rho)M^3w_1(w_1+w_2)
       \sup_{b\text{ oriented and meeting }\square}
                                                      |e^{iqZ(b)}-1|.
\end{equation}
If $q\le1$ and
$q|Z(b)|\le\frac12q^{1-\varepsilon/2}$ on these bonds,
then
\begin{equation}\label{eq:rg-step-first-vertex-small-bound}
 \|\mathcal V_0(Z,\square)\|_{h_0/2}
                 \le6(1+\rho)M^3q^{1/2-\varepsilon/2}
                 \le12M^3q^{1/2-\varepsilon/2}.
\end{equation}
\end{proposition}

\begin{proof}
At a nearest-neighbor pair the six paths reduce to
the bond, so the quotient relations read
\[
 f(x)+\chi(x,x^{(3)})=u_{\mathscr T}(c,c_3)f(x^{(3)}),
 \qquad
 \bar f(x)+\bar\chi(x,x^{(3)})
                         =u_{\mathscr T}(c,c_3)^{-1}\bar f(x^{(3)}).
\]
The first half of the summed vertex is therefore
one half of the integral of
\[
 \bar f(x)^TM_3(x)(e^{iqZ(c,c_3)}-1)
                                  u_{\mathscr T}(c,c_3)f(x^{(3)}).
\]
Reindexing the second half by the measure-preserving
involution $x\mapsto x^{(3)}$ gives the same
integrand: both $Z$ and $\mathscr T$ change sign
on the reversed bond. Only the type-$3$ hop
depends on the connection, and
$u_{\mathscr T+Z}=e^{iqZ}u_{\mathscr T}$.
The difference of the two Dirac operators is
thus the negative of this hopping coefficient,
proving \eqref{eq:rg-step-first-vertex-evaluation}.
Every product contains two antiperiodic fields,
so the reindexing is legitimate on the torus.

For lattice symmetry use
$g^{-1}(x^{(3)})=(g^{-1}x)^{(3)}$ and
$S^{-1}M_3(x)S=M_3(g^{-1}x)$ for the spin
lift $S$. Each half transforms separately, and
reindexing maps $g\square\times O_{\mathrm{cub}}$
onto $\square\times O_{\mathrm{cub}}$ with the
same measure. For a real gauge, the two generators
in each term have the same first photon site and
opposite bar indices, so their phases cancel.
Charge conjugation interchanges the two halves.
Its algebraic rule on the first half gives
\[
 \mathrm{Cg}[\bar\psi(x)^Tm(\psi(x)+\chi(x,x^{(3)}))]
       =(\bar\psi(x)+\bar\chi(x,x^{(3)}))^T
                              (\mathsf C^{-1}m^T\mathsf C)\psi(x).
\]
The sign from exchanging the odd factors cancels
the minus sign in the image of $\bar\psi$.
As
\[
 \mathsf C^{-1}M_3(x)^T\mathsf C=M_3(x^{(3)}),
\]
the first half at $-Z$ becomes the second at $Z$;
the reverse relation gives the other half.
Finally
\[
 \varkappa^T\overline{M_3(x)}\varkappa=M_3(x),\qquad
 \overline{e^{\pm iq(-\overline Z)}-1}=e^{\pm iqZ}-1.
\]
The antilinear map consequently preserves each
half in its reality identity. This proves
\eqref{eq:rg-step-first-vertex-symmetries}.

Put $\delta_x=e^{iqZ(c,c_3)}-1$ and
$\delta'_x=e^{-iqZ(c,c_3)}-1$. Collecting the
ordinary bilinear in \eqref{eq:rg-step-first-vertex},
its coefficient is
$[M_3(x)\delta_x+M_3(x^{(3)})\delta'_x]/96$.
The other two products have coefficients
$M_3(x)\delta_x/96$ and
$M_3(x^{(3)})\delta'_x/96$, and weights
$w_1w_2$.
The ordinary product has weight $w_1^2$.
The chamber sites in the cube and the ordered
pairs with these first sites are distinct, so
these are exactly the coefficients, up to
ordering signs. Each endpoint matrix satisfies
$\sum_{\beta,\beta'}|M_3(\cdot)_{\beta\beta'}|
\le12(1+\rho)$, since
$M_3=3(\rho\mp\gamma_\mu)$ and $\gamma_\mu$
has four entries of modulus one.
There are $48M^3$ chamber sites in the cube.
The triangle inequality for the collected
bilinear therefore gives
\eqref{eq:rg-step-first-vertex-bound}.

Under the smallness condition,
$|e^{iqZ(b)}-1|\le2q|Z(b)|\le q^{1-\varepsilon/2}$.
For $w=h_0/2$,
$w_1(w_1+w_2)=\frac14q^{-1/2}(1+q^\varepsilon)
\le\frac12q^{-1/2}$. Inserting these bounds proves
\eqref{eq:rg-step-first-vertex-small-bound}.
\end{proof}

For $Z=\mathscr Z$ and the evaluation at
$\Psi^f,\bar\Psi^f$, the first term of
\eqref{eq:rg-step-first-exponent} is the image of
the summed vertex. Under the additional charge
condition $96M^3q^{1/4}\le1$, its bound becomes
$\|\mathcal V_0(\mathscr Z,\square)\|_{h_0/2}
\le q^{1/4-5\varepsilon}$ whenever the displayed
bond smallness condition holds. These statements
apply before any expansion of the finite integral.

The first-step density identity retains every
vacuum and determinant factor:
\begin{equation}\label{eq:rg-step-first-density}
 \begin{split}
 (\mathbf T_0^x\mathscr F_0^x[\varepsilon^{\mathrm v},m,0])
                                    (A',\mathscr C';\Psi')\\
  ={}&\det D_1(0)^{-1}
       \exp\Big\{-\langle\bar\Psi',D_1(\mathscr T')\Psi'\rangle_1
                   -Lm\langle\bar\psi',\psi'\rangle\\
      &\hspace{14mm}
                   -L^3(\varepsilon^{\mathrm v}
                        +\delta\varepsilon_0^0)n_1
                   +\sum_X\mathfrak P_0(X,\mathscr C)\Big\}\\
       &\times
              \frac{\det F_0(\mathscr T)}{\det F_0(0)}
        \int\dd\mu_I(W)\,\dd\widehat\mu_0(\zeta)\,e^{V_0^x}.
 \end{split}
\end{equation}
The two determinants of $F_0$ have size $192n_0$.
The original level-zero density, integrated before
the block transformation, is
\begin{equation}\label{eq:rg-step-initial-integrated-density}
 \int\mathrm D\Psi\,\mathscr F_0^x
                     [\varepsilon^{\mathrm v},m,0](A,C;\Psi)
     =e^{-\varepsilon^{\mathrm v}n_0+\sum_X\mathfrak P_0(X,C)}
               \frac{\det(D_{A+C}+\mu+m)}{\det D_0(0)}.
\end{equation}
This follows directly from the weighted determinant
identity. The numerator need not be invertible;
the nonzero denominator uses $\mu>0$.


\section{Localizing the fluctuation activity}\label{sec:rg-nonlinear-step}

The finite fluctuation integral in
\eqref{eq:rg-step-fluctuation-exponent} has a logarithm which is a
sum of local activities. Its construction has two localizations.
The first resolves the change of background and reconstructs
the coarse mean; the second resolves the fermion fluctuations.
After these finite operations, a connected-polymer expansion
gives the logarithm. The estimates below apply also to the first
step, where the fluctuation generators have chamber measure
$1/48$.

\subsection{Contour domains and localization}
\label{subsec:rgnl-domains}

Keep the admissible sides, mass and charge of
\cref{sec:rg-step}, and assume in addition
$\ell\ge Mr_k$. Put $e=e_k$, $h=(h_{k,1},h_{k,2})$.
The weakened photon and fermion primitives are taken on the
polydisc
\[
 \mathrm{Pd}_k=\{s=(s',s''):|s_{\mathrm{pc}}|\le M^{\alpha_0}\}.
\]
Here $s'=(s'_M,s'_{LM})$ contains the photon reconstruction
and square-root parameters, of kinds $\pi_M,\pi_{LM}$ in
\exteq{37}{MaxwellBlockManuscript}; $s''$ contains the
fermion $M$-cube parameters $\pi_k$ of
Section~\extref{4.1}{FermionBlockManuscript}.
We assume their analytic kernel and H\"older estimates,
component restriction, row vanishing and coefficient support
with common constants $(\alpha_0,\gamma_w,\kappa_w;C_w)$,
where $\kappa_w=\kappa+\kappa_0+3$.
The photon hypotheses are those of
Proposition~\extref{5.1}{MaxwellBlockManuscript},
\exteq{42}{MaxwellBlockManuscript} and the component
restriction following them. The fermion hypotheses are those of
Propositions~\extref{4.1}{FermionBlockManuscript} and
\extref{4.3}{FermionBlockManuscript}, including
\exteq{57}{FermionBlockManuscript},
\exteq{58}{FermionBlockManuscript}, and their covariance
identities. At parameter value one all families equal the
exact operators. Their gauge, lattice, transpose and reality
actions are the actions of those exact operators with the
parameters transformed accordingly. At level zero we retain
the stated diagonal families and the chamber covariance
with measure $1/48$.

The quantitative assumptions are those of \cref{thm:rgnl-current-step}.

For $\zeta\in\mathrm{Box}_k$, write
\begin{equation}\label{eq:rgnl-weakened-photon}
 Z(s')=H_k^{\mathrm{ph}}(s')C C_k^{1/2}(s')\zeta .
\end{equation}
In this formula $C$ is the real photon coordinate map of
\exteq{21}{MaxwellBlockManuscript}. The diagonal photon
family is used at $k=0$. For a bond field $Z$ set
\[
 |Z|_*=\max\{\|Z\|_\infty,\|\partial Z\|_\infty,
                              \|\delta_\alpha\partial Z\|_\infty\}.
\]
The three thresholds defining $R_k$ are
$\theta_i=e^{-3/4+i\varepsilon}$, $i=1,2,3$.
Thus $\widetilde R_k^-$ consists of the permitted real-gauge
part plus a field satisfying the three strict bounds
$(1-e^\varepsilon)\theta_i$.
Introduce the radii
\begin{equation}\label{eq:rgnl-contour-radii}
 r_1=\tfrac12e^{-3/4+9\varepsilon/2},\qquad
 r_2=e^{-3/4+7\varepsilon},\qquad
 r_4=e^{-1/4+3\varepsilon}.
\end{equation}

\begin{proposition}\label{prop:rgnl-contour-domains}
For every $s\in\mathrm{Pd}_k$ and
$\zeta\in\mathrm{Box}_k$,
\[
 |Z(s')|_*\le\tfrac12e^{-\varepsilon/2}.
\]
The radii satisfy $2\le r_2,r_4\le r_1$ and
$(r_i-1)^{-1}\le2/r_i$, $i=2,4$.
If $A\in\widetilde R_k^-$, $\mathscr C\in\mathfrak X_k$,
and $|Z|_*\le\tfrac12e^{-\varepsilon/2}$, then
\[
 A+tZ\in\widetilde R_k,\qquad
 A+\mathscr C+tZ\in\widetilde{\mathcal G}_k
                  \quad (|t|\le r_1).
\]
\end{proposition}

\begin{proof}
The square-root kernel is bounded by
$C_w\exp(-\gamma_w d(\Upsilon,\Upsilon'))$.
The summability constant $c_\gamma$ and
the coordinate bound $|\zeta(\Upsilon')|\le p_{0,k}
\le e^{-\varepsilon/4}$ give
\[
 \sup_\Upsilon|(C_k^{1/2}(s')\zeta)(\Upsilon)|
                                      \le C_wc_\gamma e^{-\varepsilon/4}.
\]
The coordinate map is real. Its bound
$\|Cv\|_\infty\le c_C\sup_\Upsilon|v(\Upsilon)|$
also gives the absolute row bound
$\sum_\Upsilon|C(b,\Upsilon)|\le c_C$:
apply the real bound to the signs of the entries of a fixed
row. Therefore
\[
 \|CC_k^{1/2}(s')\zeta\|_\infty
                                  \le c_CC_wc_\gamma e^{-\varepsilon/4}.
\]
Each of the reconstruction, derivative and H\"older derivative
kernels has the same summable exponential bound. Consequently
\[
 |Z(s')|_*\le c_CC_w^2c_\gamma^2 e^{-\varepsilon/4}
                                          \le\tfrac12 e^{-\varepsilon/2},
\]
where the last inequality is
$2c_CC_w^2c_\gamma^2e^{\varepsilon/4}\le1$.
At level zero the polynomial diagonal family has the same
row bounds: its constant is
$4\exp(2\gamma_w)M^{27\alpha_0}\le C_w$.

Since $\varepsilon\le1/200$,
\[
 \frac{r_2}{r_1}=2e^{5\varepsilon/2}\le1,\qquad
 \frac{r_4}{r_1}=2e^{1/2-3\varepsilon/2}
                                      \le2e^{5\varepsilon/2}\le1.
\]
The small-charge condition
$88M^3e^{1/4-6\varepsilon}\le1$ gives
$r_4^{-1}=e^{1/4-3\varepsilon}\le1/88$ and
$r_2^{-1}\le r_4^{-1}$.
In particular both radii are at least $2$, and their stated
inverse bounds follow.

For $|t|\le r_1$ the three norms of $tZ$ are at most
$\frac14e^{-3/4+4\varepsilon}$, strictly less than
$e^\varepsilon\theta_3$. Adding this field fills less than the
margin in each of the three defining bounds of
$\widetilde R_k^-$. The permitted gauge part is unchanged.
Hence $A+tZ\in\widetilde R_k$.
Finally $\widetilde R_k+\mathfrak X_k
\subset\widetilde{\mathcal G}_k$ gives the second inclusion.
\end{proof}
The first-step vertex estimate
\eqref{eq:rg-step-first-vertex-small-bound} now applies to
$Z(s')$ for every weakening parameter. In particular,
$96M^3e^{1/4}\le1$ gives its bound
$e^{1/4-5\varepsilon}$ in the $h/2$ norm.

Let $B=A+\mathscr C$.
The arguments of $H_k^\#$ are independent pairs
$(\Phi,\bar\Phi)$; in $H_k^\#\Phi$ the single symbol denotes
this pair. Write
$\psi^\#=(\psi,\chi;\bar\psi,\bar\chi)$ for the ordinary,
pair, barred ordinary and barred pair generators, respectively.
Then
\[
 \begin{aligned}
 H_k^\#(s'',B)\Phi
   ={}&(H_k(s'',B)\Phi,\delta_{\alpha,B}H_k(s'',B)\Phi;\\
       &
       \bar H_k(s'',B)\bar\Phi,
                       \delta_{\alpha,-B}\bar H_k(s'',B)\bar\Phi),\\
 \bar H_k(s'',B)&=\mathsf C^{-1}H_k(s'',-B)\mathsf C.
 \end{aligned}
\]
Define
$J_k^\#(s'',B,Z)=H_k^\#(s'',B+Z)-H_k^\#(s'',B)$ and
\[
 \widetilde\Gamma_k(s'',B)(W,\bar W)
                         =(\Gamma_k(s'',B)W,\bar W),\qquad
 W_k^\#(s'',B)=H_k^\#(s'',B)
                         \big(\Gamma_k(s'',B)W,\bar W\big).
\]
The notation $J_k,\bar J_k$ denotes the components without
H\"older quotient; $\mathrm{Wc},\overline{\mathrm{Wc}}$ denotes
the corresponding components of $W_k^\#$.
The auxiliary generators $W$ have weight
$\mathsf w_W=1$ for $k\ge1$ and $\sqrt{48}$ for $k=0$.

These fields are continuous in all closed polydisc and box
variables and analytic in the open parameter polydisc and the
connection variables. For the H\"older rows this follows
directly from
\[
 (\delta_{\alpha,B}H_k)((x,x'),u)
   =d_X(x,x')^{-\alpha}
       \big(U_B(x,x')H_k(x',u)-H_k(x,u)\big).
\]
The transport is entire in $B$. \Cref{prop:rgnl-contour-domains}
places every connection on the closed contours inside the
open domain of the primitives. Thus compositions at
$B+tZ(s')$ have the stated regularity also when the contour
variable is on its circle.

In the following localization, cubes and polymers are sets
of photon sites. Containment refers to all their lattice sites,
and two cubes meet when they share a lattice site.
Let $P_0(X)$ be the parameter cubes contained in $X$.
For a set $P$ of parameter cubes disjoint from $P_0(X)$, let $Y'(X,P)$ be the union of
the $M$-cubes meeting $X$ or a cube of $P_0(X)\cup P$.
For a continuous analytic family $G(s,X)$ define
\begin{align}
 \mathcal T_PG(X)
  &=\int_{[0,1]^P}\dd s_P\,
       \partial_PG(s,X)
       \big|_{s=1\text{ on }P_0(X),\,
                    s=0\text{ outside }P_0(X)\cup P},
                          \label{eq:rgnl-interpolation}\\
 \operatorname{Loc}G(Y)
  &=\sum_{(X,P):Y'(X,P)=Y}\mathcal T_PG(X).
                          \label{eq:rgnl-localization}
\end{align}
Repeated one-variable integration gives the useful identity
\begin{equation}\label{eq:rgnl-finite-difference}
 \mathcal T_PG(X)=
    \sum_{Q\subseteq P}(-1)^{|P|-|Q|}
                                 G(\mathbf1_{P_0(X)\cup Q},X).
\end{equation}
In particular, coefficient locality can be checked at
indicator families; no statement about the convergence of
the later logarithm is used here.

\begin{lemma}\label{lem:rgnl-resummation}
Suppose
$\|G(s,X)\|\le B_0\exp(-\kappa'd_M(X))$ on the parameter
polydisc, where $\kappa_0\le\kappa'\le\kappa$.
Assume $\mathcal T_PG(X)=0$ whenever $Y'(X,P)$ is
disconnected. Then
\begin{gather}
 \sum_Y\operatorname{Loc}G(Y)=\sum_XG(1,X),
                         \label{eq:rgnl-localization-sum}\\
 \|\operatorname{Loc}G(Y)\|
   \le4c_XK_{\mathrm{pol}}B_0
             \exp\big(-(\kappa'-\kappa_0-1)d_M(Y)\big).
                         \label{eq:rgnl-resummation-bound}
\end{gather}
\end{lemma}

\begin{proof}
Applying $f(1)=f(0)+\int_0^1f'(t)\dd t$ in each parameter
outside $P_0(X)$ gives $G(1,X)=\sum_P\mathcal T_PG(X)$.
The sums are finite. After removing the vanishing disconnected
terms, each term belongs to exactly one polymer $Y'(X,P)$.
This proves \eqref{eq:rgnl-localization-sum}.

Put $R=M^{\alpha_0}$. Cauchy's estimate on a polydisc of
radius less than $R-1$ about each integration point, followed
by the limit of these radii to $R-1$, gives
\[
 \|\mathcal T_PG(X)\|
       \le B_0\exp(-\kappa'd_M(X))(R-1)^{-|P|}.
\]
If $Y'(X,P)=Y$, then $X\subseteq Y$ and
$|Y\setminus X|_M\le N_c|P|$, since every added $M$-cube
meets a parameter cube and each parameter cube meets at
most $N_c=(2L+1)^3$ such cubes.
To compare their tree lengths, order the cubes of $Y\setminus X$
by graph distance from $X$. Each added cube is adjacent
to a preceding one. Extend a net by a segment from
a point of that preceding closed cube to their common
face, of length at most $M$. Starting from an
arbitrarily near-minimal net for $X$ gives
$d_M(Y)\le d_M(X)+|Y\setminus X|_M$.
Conversely every net for $Y$ is a net for $X$.
Thus $0\le d_M(Y)-d_M(X)\le N_c|P|$.
Writing $a_X=\log(6c_X)$, the condition
$R-1\ge\exp(N_c\kappa+a_X)$ implies
\[
 \|\mathcal T_PG(X)\|
 \le B_0\exp(-(\kappa'-\kappa_0)d_M(Y))
             \exp(-\kappa_0d_M(X))\exp(-a_X|P|).
\]
At most $3|Y|_M$ parameter cubes are contained in $Y$.
Hence
\[
 \sum_P\exp(-a_X|P|)
       \le(1+(6c_X)^{-1})^{3|Y|_M}
       \le \exp((d_M(Y)+1)/2).
\]
The polymer counting and size estimates give
\[
 \sum_{X\subseteq Y}\exp(-\kappa_0d_M(X))
   \le K_{\mathrm{pol}}|Y|_M
   \le2c_XK_{\mathrm{pol}}\exp(d_M(Y)/2).
\]
The product of the last two bounds is at most
$4c_XK_{\mathrm{pol}}\exp(d_M(Y))$. Summing the estimate for
$\mathcal T_PG(X)$ proves
\eqref{eq:rgnl-resummation-bound}.
\end{proof}

\subsection{The two finite localizations}
\label{subsec:rgnl-finite-localizations}

For a function analytic near the closed disk $|t|\le r$, $r>1$,
write
\[
 \mathcal C_rf=\frac1{2\pi i}\int_{|t|=r}
                                      \frac{f(t)}{t(t-1)}\,\dd t .
\]
This is $f(1)-f(0)$ and satisfies
$\|\mathcal C_rf\|\le(r-1)^{-1}\sup_{|t|=r}\|f(t)\|$.
Let $(m,\mathscr E)\in\operatorname{Ball}_k^x(4)$ and set
$m'=m+m(\mathscr E)$, $\mathscr E^r=R\mathscr E$.
For $k\ge1$ define
\begin{align}
 G_1(s,X)&=\mathcal C_{r_1}
                \langle\bar\Psi,M_k(X,B+tZ(s'))\Psi\rangle_1,
                                      \label{eq:rgnl-G1}\\
 G_2(s,X)&=\mathcal C_{r_2}\,
       \mathscr E^r(X;A+tZ(s'),\mathscr C,
                    \psi^\#+tJ_k^\#(s'',B,Z(s'))\Psi),
                                      \label{eq:rgnl-G2}\\
 G_3(s,\square)&=\mathcal C_{r_2}(-m')
       \langle\bar\psi+t\bar J_k(s'',B,Z(s'))\bar\Psi,
                         \psi+tJ_k(s'',B,Z(s'))\Psi\rangle_\square .
                                      \label{eq:rgnl-G3}
\end{align}
The symbols $M_k(X,\cdot)$ denote the local Dirac primitives.
A function defined for one cube is zero on the other polymers.
Put
\[
 E_i^*(Y)=\operatorname{Loc}G_i(Y),\qquad
 E_i(Y)=\operatorname{sub}_{Y,B}E_i^*(Y),\qquad
 E_{\le3}=\sum_{i=1}^3E_i.
\]
The substitution is \eqref{eq:rg-step-local-substitution}.
At $k=0$ there is no $\Psi$ and no substitution:
$G_1(s,\square)=\mathcal V_0(Z(s'),\square)$,
$G_2=G_3=0$, and $E_{\le3}=\operatorname{Loc}G_1$.
All parameter kinds occur in this first localization.

For the second localization, which uses only $\pi_k$, let
$v=W_k^\#(s'',B)$ and define
\begin{align}
 G_{\le3}^{\,W}(s'',X)&=E_{\le3}(X;A,\mathscr C,\zeta,
                                                      \psi^\#+v),
                                      \label{eq:rgnl-GW}\\
 G_4(s'',X)&=\mathcal C_{r_4}
                  \mathscr E^r(X;A,\mathscr C,\psi^\#+tv),\\
 G_5(s'',\square)&=\mathcal C_{r_4}(-m')
                 \langle\bar\psi+t\overline{\mathrm{Wc}},
                                   \psi+t\mathrm{Wc}\rangle_\square,\\
 F(Y)&=\operatorname{Loc}(G_{\le3}^{\,W}+G_4+G_5)(Y).
                                      \label{eq:rgnl-finite-activity}
\end{align}
At $k=0$ the term $G_4$ is zero and $m'=m$.

\begin{proposition}\label{prop:rgnl-finite-locality}
All finite expressions just defined exist on
$\widetilde R_k^-\times\mathfrak X_k\times\mathrm{Box}_k$.
The element $F(Y)$ is even, has ordinary legs over $Y$ and
pair legs admissible for $Y$, and has fluctuation generators
only at unit indices in $Y$. Its coefficients read $A,\mathscr C$
only on bonds meeting $Y$ and read $\zeta$ only through the
coordinates $\zeta(\Upsilon)$ with $\supp\Upsilon\subset Y$.
They are continuous in these arguments and analytic in the
connection variables. The same assertions hold for $E_{\le3}(Y)$,
which has no fluctuation generators.

The disconnected terms in both localizations vanish, and
\begin{equation}\label{eq:rgnl-two-localization-sums}
 \sum_YE_i^*(Y)=\sum_XG_i(1,X),\qquad
 \sum_YF(Y)=\sum_X(G_{\le3}^{\,W}+G_4+G_5)(1,X).
\end{equation}
The finite expressions on $Y$ depend on the datum only through
$m'$ and the $\mathscr E^r(X)$ with $X\subseteq Y$.
\end{proposition}

\begin{proof}
We give the component argument which is responsible for
locality. Let $\mathrm{Ps}=P_0(X)\cup P$, let
$Y'=Y'(X,P)$ and let $Y'_1$ be its component containing $X$.
Set the weakening parameters to zero outside $\mathrm{Ps}$.
For each photon parameter kind, take components linked by shared
lattice sites; for the fermion kind $\pi_k$, take face-connected
components of the active partition $M$-cubes. In either case,
a component $\mathrm{Lnk}$ of $\mathrm{Ps}$ is contained in the
polymer formed by the $M$-cubes meeting it. This polymer lies in one
component of $Y'$. If a site of $Y'_1$ lies in $\mathrm{Lnk}$,
or in the enlarged parameter cube of a photon kind, that
polymer has a cube in common with $Y'_1$, and hence lies
entirely in $Y'_1$.

The row-vanishing property of the weakened photon families
and their component restriction give, on any fine bond $b$
meeting $Y'_1$,
\begin{equation}\label{eq:rgnl-restricted-photon}
 Z(s')(b)=
  \sum_{\substack{b',\Upsilon,\Upsilon'\subset Y'_1}}
   H_k^{\mathrm{ph}}(s')(b,b')C(b',\Upsilon)
          C_k^{1/2}(s')(\Upsilon,\Upsilon')\zeta(\Upsilon').
\end{equation}
Indeed, the reconstruction row is zero unless $b$ lies in
one link-component; its columns lie in that component.
The coordinate map has columns supported in the $L$-blocks
of the ends of $b'$, and each block lies in the $M$-cube of
the corresponding end. The square-root row then has columns
in its own link-component. Each component involved is
attached to $Y'_1$ by the preceding argument.
These rows are unchanged if the parameters outside $Y'_1$
are changed while their common zero set outside $\mathrm{Ps}$
is kept.

For a chamber site $x=(c,g)$ at distance at most $1$ from
$Y'_1$, put $K_c=\{c':|c'-c|_\infty\le2L\}$, as in
\cite[Proposition~4.1]{FermionBlockManuscript}.
The fermion row is zero unless every parameter cube
meeting $K_c$ is present. These partition cubes form a
face-connected set. When the row is nonzero, their active
face-connected component contains $c$ and a site of $Y'_1$,
so it lies in $Y'_1$.
The row of $H_k(s'',B)$, and its barred counterpart,
then has columns in that component and reads $B$ only on
bonds meeting it. At a column of this row, the corresponding
row of $\Gamma_k(s'',B)$ has columns in the same component.
These conclusions also hold at level zero: the diagonal
factor is
\[
 \prod_{\mathrm{pc}\in\pi_0:\mathrm{pc}\cap K_c\ne\varnothing}
                                                       s''_{\mathrm{pc}},
\]
and the kernel of that multiplication operator is
$48\mathbf1_{x=x'}$ times this factor. Its composition with
$\Gamma_0$ uses the measure $1/48$.

An endpoint of a pair admissible for $X$ is at distance
at most $\eta_k\le1$ from $X$. Its transport reads only
bonds contained in $X$, or the one type-$3$ bond meeting $X$.
Thus the preceding row argument applies to every row
used in $J_k^\#$ and $W_k^\#$.
The H\"older row has only the two columns $x,x'$ and
the entire transport coefficient displayed above.
Finite compositions preserve this dependence: the columns
of $O_2O_1$ are the union of the columns of $O_1$ at the
finitely many columns of $O_2$, and its read bonds are the
union of their read bonds. Entries of the transport and
reconstruction rows are entire in the connections, and finite
compositions of these algebraic rows retain this property.
Compositions involving weakened primitives are holomorphic
on their stated primitive domains. These are the finite-row
properties used here; the weakened operators themselves
are not asserted to have finite range.

For $G_1$ at $k\ge1$, the local primitive $M_k(X,\cdot)$
has both indices in $X$ and reads only bonds meeting $X$.
For $G_2,G_3$, the incoming activity or mass form has ordinary
legs in $X$ and admissible pair legs; the substitution by
$J_k^\#\Psi$ introduces $\Psi$ only in $Y'_1$.
At $k=0$ the vertex has precisely these ordinary and
type-$3$ pair legs and reads $Z$ only on bonds meeting
its cube. Equation~\eqref{eq:rgnl-restricted-photon}
therefore proves the same component dependence for every
first-localization function.

If $Y'$ is disconnected, choose an $M$-cube outside $Y'_1$.
It cannot meet $X$ or a parameter cube contained in $X$,
so it meets a cube of $P$ outside $Y'_1$. The integrand is
independent of that parameter. Pairing the subsets in
\eqref{eq:rgnl-finite-difference} according to whether they
contain this cube proves that the term is zero.
For a term assigned to a connected $Y$, all parameter cubes
in $P_0(X)\cup P$ and $X$ lie in $Y$.
Equation~\eqref{eq:rgnl-finite-difference} now gives the
claimed locality of $E_i^*$.

The reconstruction substitution adds only ordinary legs
over $Y$ and boundary pairs $(x,x^{(3)})$ admissible for
$Y$. Its finite rows at a unit site $y$ read bonds contained
in $\Delta_y$ or type-$3$ bonds from this block. These bonds
meet $Y$ since $\Delta_y\subset Y$. The barred rows have
the same support. Thus it preserves the claimed coefficient
locality and removes $\Psi$.
For the second localization, the row argument for
$H_k^\#\widetilde\Gamma_kW$ gives fluctuation generators
only in $Y'_1$. It applies to each of the three functions
in \eqref{eq:rgnl-finite-activity}. The same disconnected
parameter cancellation proves the second assertion.
\Cref{lem:rgnl-resummation} gives both sum identities.

Regularity follows from the established contour domains,
the primitive regularity, and finite sums and substitutions.
Contour integrals preserve analyticity: their Riemann sums
converge uniformly on compact sets. The same argument
applies to finite parameter integrals. The incoming
activities and mass forms are even, and each substitution
sends generators to odd linear combinations, so every
expression remains even.
Finally $G_1$ is independent of the datum; $G_2,G_4$
use $\mathscr E^r(X)$ and $G_3,G_5$ use $m'$.
Tracing the two localizations shows that no incoming
polymer outside $Y$ enters the result.
\end{proof}

The proof identifies the primitive entries determining a
finite expression on $Y$. With parameters zero outside $Y$,
the photon field reads only
$H_k^{\mathrm{ph}}(s')(b,b')$ for $b$ meeting $Y$,
$b'\subset Y$, and
$C_k^{1/2}(s')(\Upsilon,\Upsilon')$ with both supports
in $Y$. Its remaining factor is the fixed coordinate map.
The fermion fields use $H_k,\bar H_k$ at chamber sites
over $Y$, with unit columns in $Y$, at $B+Z(s')$ and
$B$, and $\Gamma_k(s'',B)$ with both unit indices in $Y$.
At zero, the kernel sums have measure $1/48$ and the
diagonal rows have the factor $48$ already recorded.
The Dirac primitive has both indices in $X\subseteq Y$
and connection $B+tZ(s')$ on $|t|=r_1$.
H\"older transports read only bonds meeting $Y$.
The reconstruction substitution contains only the fixed
finite $T_k^{(i)},\bar T_k^{(i)}$ rows at unit sites of its
own subpolymer.

Consequently two systems with the stated hypotheses on this torus,
with the same $m'$ and the same incoming $\mathscr E^r(X)$
for every $X\subseteq Y$, give the same finite expressions on $Y$
if these entries agree
for every parameter family zero outside $Y$, at the
connections formed with their respective photon fields.
The restricted photon sum first gives equality of the
fields on bonds meeting $Y$; the transports then agree. The
incoming activities agree by the common-data hypothesis and
their locality.
The listed fermion and Dirac entries give equal first
functions, their localizations and reconstruction
substitutions. They give equal second functions and
localizations in turn, and then equal activities $K(Y)$ and
$H(Y)$ in \eqref{eq:rgnl-Mayer-activity} and
\eqref{eq:rgnl-integrated-activity}, and each term of the
connected series \eqref{eq:rgnl-log-activity}.
Agreement of the photon
entries and agreement of the fermion entries as functions
on the full operator domain is sufficient: if the two
composed connections differ outside $Y$, local row
dependence removes that difference first.
The same comparison passes to the convergent logarithm
once its convergence has been established.

\subsection{Bounds before fluctuation integration}
\label{subsec:rgnl-finite-bounds}

Set
\[
 \begin{gathered}
 \kappa_w=\kappa+\kappa_0+3,\qquad
 a_1=\kappa-\kappa_0-1,\qquad a_2=\kappa-2\kappa_0-2,\\
 C_W=28C_w^2c_\gamma,
 \qquad H_{\mathrm{cl}}=36c_X^2K_{\mathrm{pol}}^2e^{1/4-5\varepsilon}.
 \end{gathered}
\]
The notation $H_{\mathrm{cl}}$ is a scalar cluster bound. For a
scalar-component kernel $O(a,b)$ with column measure $\nu(b)$, write
$|O|_{1,\infty}=\sup_a\sum_b\nu(b)|O(a,b)|$, including spin indices.
This is the row norm of \cite[Theorem~B]{NormalizedContraction}.
The finite primitive bounds imply
\begin{align}
 \|\langle\bar\Psi,M_k(X,B)\Psi\rangle_1\|_{w_3}
 &\le4c_Xc_\gamma C_wM^3w_3^2
                   \exp(-(\kappa_w-1)d_M(X)),\quad k\ge1,
                                      \label{eq:rgnl-M-bound}\\
 |H_k^\#(s'',B)|_{1,\infty}&\le4C_w,\\
 |J_k^\#(s'',B,Z)|_{1,\infty}
                  &\le16C_we^{3/4-9\varepsilon/2},\\
 \mathsf w_W|H_k^\#(s'',B)\widetilde\Gamma_k(s'',B)|_{1,\infty}
                  &\le C_W.           \label{eq:rgnl-W-bound}
\end{align}
Here $B=A+\mathscr C$ and $|Z|_*\le\frac12e^{-\varepsilon/2}$.
To see the first bound, count $4M^3|X|_M$ rows and sum
the exponential kernel in its second index, giving
$c_\gamma C_w \exp(-\kappa_wd_M(X))$ per row.
Then $|X|_M\le c_X\exp(d_M(X))$ gives
\eqref{eq:rgnl-M-bound}.
For $H_k^\#$ the primitive bound controls the ordinary
and short-pair rows by $C_w$. A long pair, $d_X\ge1$,
has its H\"older quotient bounded by four times the
ordinary sup norm. At level zero every pair is long.
The monomial matrices $\mathsf C^{\pm1}$ preserve row norms,
so the barred rows have the same bounds.
Apply $\mathcal C_{r_1}$ to
$H_k^\#(s'',B+tZ)$ to obtain the bound for $J_k^\#$.
Finally $|\Gamma_k|_{1,\infty}\le C_wc_\gamma$ by kernel
summation, using $1/48$ at level zero. Submultiplicativity
and $\mathsf w_W\le7$ give \eqref{eq:rgnl-W-bound}.

\begin{proposition}\label{prop:rgnl-finite-bounds}
For the data of \cref{prop:rgnl-finite-locality},
\begin{align}
 \|E_{\le3}(Y)\|_{h/2}
       &\le8c_XK_{\mathrm{pol}}e^{1/4-5\varepsilon}\exp(-a_1d_M(Y)),
                                      \label{eq:rgnl-E3-bound}\\
 \|F(Y)\|_{(h/4;\mathsf w_W)}
       &\le H_{\mathrm{cl}}\exp(-a_2d_M(Y)).
                                      \label{eq:rgnl-F-bound}
\end{align}
The element $F(Y)$ is the sum of a term independent of
$(m,\mathscr E)$ and a term complex linear in these data.
\end{proposition}

\begin{proof}
Extraction gives
\[
 \|\mathscr E^r(X)\|_h\le2\|\mathscr E\|_k^x\exp(-\kappa d_M(X))
                   \le12e^{1/4-7\varepsilon}\exp(-\kappa d_M(X)),
 \qquad |m'|\le8e^{3/4-8\varepsilon}.
\]
The mass bound uses
$c_{\mathrm{fl}}e^\varepsilon\le1$.
With chamber measure $\eta_k^3/48$, direct counting gives
\[
 \|\langle\bar\psi,\psi\rangle_\square\|_{(w_1,w_2)}
                                                    =4M^3w_1^2.
\]
The number of monomials is
$4\cdot48\eta_k^{-3}M^3$ and their common coefficient
is $\eta_k^3/48$.

In the first localization give $\Psi$ the weight
$w_3=C_Th_{k,1}$ and the old generators the weights $h/2$.
The contour bound and \eqref{eq:rgnl-M-bound} give
\[
 \|G_1(s,X)\|\le
 16c_Xc_\gamma C_wC_T^2M^3e^{1/4-9\varepsilon/2}
                 \exp(-(\kappa_w-1)d_M(X))
 \le e^{1/4-5\varepsilon}\exp(-\kappa d_M(X)).
\]
The last step uses
\[
 16c_Xc_\gamma C_wC_T^2M^3e^{\varepsilon/2}\le1.
\]
At $k=0$ the weighted vertex instead gives
$12M^3e^{1/2-\varepsilon/2}\le e^{1/4-5\varepsilon}$.
Moreover,
\[
 r_2|J_k^\#|_{1,\infty}w_3
     \le16C_wC_Te^{3\varepsilon/2}h_{k,2}
     \le\tfrac12h_{k,2}.
\]
The substitution bound therefore takes the integrand norm
of $G_2$ to the $h$ norm of $\mathscr E^r$, and that of
$G_3$ to the $h$ norm of the mass form. It gives
\[
 \|G_2(s,X)\|\le24e^{1-14\varepsilon}\exp(-\kappa d_M(X)),
 \qquad \|G_3(s,\square)\|\le64M^3e^{1-15\varepsilon}.
\]
The condition $88M^3e^{1/4-6\varepsilon}\le1$ implies
\[
 24e^{3/4-9\varepsilon}
                   +64M^3e^{3/4-10\varepsilon}\le1.
\]
Thus $\sum_{i\le3}G_i$ is bounded by
$2e^{1/4-5\varepsilon}\exp(-\kappa d_M(X))$.
Apply \cref{lem:rgnl-resummation}, using the disconnected
cancellation in \cref{prop:rgnl-finite-locality}.
The reconstruction substitution costs no norm, since
$C_T(h_{k,1}+h_{k,2})/2\le w_3$.
This proves \eqref{eq:rgnl-E3-bound}.

For the second localization, the condition
$112C_w^2c_\gamma e^{2\varepsilon}\le3$ gives
\[
 r_4C_W\le\tfrac34h_{k,2},\qquad C_W\le\tfrac14h_{k,2}.
\]
Consequently addition of $W_k^\#$ maps the $h/4$ norm
to the $h/2$ norm at $t=1$ and to the $h$ norm on
the contour $|t|=r_4$. The three input bounds are
\begin{align*}
 \|G_{\le3}^{\,W}(s'',X)\|
       &\le8c_XK_{\mathrm{pol}}e^{1/4-5\varepsilon}\exp(-a_1d_M(X)),\\
 \|G_4(s'',X)\|&\le24e^{1/2-10\varepsilon}\exp(-\kappa d_M(X)),\\
 \|G_5(s'',\square)\|&\le64M^3e^{1/2-11\varepsilon}.
\end{align*}
Their sum is at most
$9c_XK_{\mathrm{pol}}e^{1/4-5\varepsilon}\exp(-a_1d_M(X))$,
again by $88M^3e^{1/4-6\varepsilon}\le1$.
Since $\kappa_0\le a_1\le\kappa$, resummation gives
\eqref{eq:rgnl-F-bound}.
At level zero $G_4=0$; all the other inequalities include
the factor $\mathsf w_W$ already, so they apply unchanged.
Linearity follows from the definitions: $G_1$ is fixed,
$G_2,G_4$ are linear in $R\mathscr E$, $G_3,G_5$ in $m'$,
and localization, the reconstruction substitution and the
contour functionals are linear.
\end{proof}

\section{The connected logarithm}
\label{subsec:rgnl-cluster-logarithm}

Two polymers are incompatible when their union is connected
through faces. In particular a polymer is incompatible with
itself. For an ordered tuple $(X_1,\ldots,X_n)$ define
$u^T(X_1)=1$ and
\[
 u^T(X_1,\ldots,X_n)=\sum_{G\text{ connected}}
                                      (-1)^{|G|},\quad n\ge2,
\]
where $G$ ranges over spanning subgraphs of the incompatibility
graph of the tuple. The coefficient is zero when that graph
is disconnected.

We will use the formal connected-graph identity
\begin{equation}\label{eq:rgnl-connected-graph-identity}
 \sum_{\{X_i\}\text{ compatible}}\prod_i z_{X_i}
 =\exp\left\{\sum_{n\ge1}\frac1{n!}
       \sum_{(X_1,\ldots,X_n)}u^T(X_1,\ldots,X_n)
                                           \prod_i z_{X_i}\right\}
\end{equation}
and the tree bound
\begin{equation}\label{eq:rgnl-tree-domination}
 |u^T(X_1,\ldots,X_n)|
       \le\#\{\text{spanning trees of the incompatibility graph}\}.
\end{equation}
Here the identity is one of formal series in commuting
variables, before substitution into a normed algebra.
For its proof, express the left side as an ordered sum with
factor $1/n!$ and multiply by
$\prod_{i<j}(1-\mathbf1_{\{X_i,X_j\}\text{ incompatible}})$.
Expanding this product sums over all graphs. A graph is
uniquely the union of its connected components. Grouping by
their vertex sets gives the exponential on the right:
an ordered list of component sizes $n_1,\ldots,n_r$ has
$n!/(n_1!\cdots n_r!)$ vertex assignments, and division
by $r!$ removes the order of its components.
Only finitely many terms contribute to a fixed degree.

For \eqref{eq:rgnl-tree-domination}, let $\mathbf G$ be
the incompatibility graph and suppose it is connected.
For a connected spanning subgraph $G$, root at $1$ and,
for $i\ne1$, choose as its parent the least labelled
neighbor at distance $d_G(i)-1$ from the root.
These edges form a tree $T(G)$ with the same root distances
as $G$. Given a spanning tree $\tau$, let $R(\tau)$ contain
its edges, all edges of $\mathbf G$ within one distance
level, and all edges from $i$ to the preceding level whose
other endpoint exceeds the parent of $i$ in $\tau$.
Then
\[
 T(G)=\tau\quad\Longleftrightarrow\quad
                                      \tau\subseteq G\subseteq R(\tau).
\]
In one direction the distances and least parents force
every edge of $G$ to belong to $R(\tau)$.
Conversely, along any edge of $R(\tau)$ the tree distance
changes by at most one. Thus a graph between $\tau$ and
$R(\tau)$ has the tree's distances, and the least neighbor
in the preceding level is its tree parent.
The connected subgraphs are therefore partitioned into
these intervals, and their signed sum is
\[
 \sum_\tau(-1)^{n-1}
           \sum_{A\subseteq R(\tau)\setminus\tau}(-1)^{|A|}
 =(-1)^{n-1}\#\{\tau:R(\tau)=\tau\}.
\]
This proves the tree bound, including the cases
of a disconnected graph and $n=1$.

For this section and the symmetry proof, write $\widehat\mu=\widehat\mu_k$ for the normalized product box law in \eqref{eq:rg-step-box-normalization}.

\begin{proposition}\label{prop:rgnl-cluster-logarithm}
Let $V(X)$ be even elements in the old and fluctuation
generators, continuous on $\mathrm{Box}_k$, with
\[
 \|V(X)\|_{(w;\mathsf w_W)}
                       \le H_{\mathrm{cl}}\exp(-a_2d_M(X)).
\]
Assume that $V(X)$ has fluctuation indices only in $X$ and
depends on $\zeta$ only through coordinates supported in $X$.
For any union $Y$ of $M$-cubes put
\begin{align}
 K(Y)&=\sum_{\substack{\{X_i\}\text{ distinct}\\\bigcup_iX_i=Y}}
                    \prod_i(\exp V(X_i)-1),\qquad K(\varnothing)=1,
                                      \label{eq:rgnl-Mayer-activity}\\
 H(Y)&=\int\dd\mu_I(W)\,\dd\widehat\mu(\zeta)\,K(Y).
                                      \label{eq:rgnl-integrated-activity}
\end{align}
For a connected $Y$ define
\begin{equation}\label{eq:rgnl-log-activity}
 E^\#(Y)=\sum_{n\ge1}\frac1{n!}
   \sum_{\bigcup_iY_i=Y}u^T(Y_1,\ldots,Y_n)\prod_iH(Y_i).
\end{equation}
The last series converges absolutely in the $w$ norm and
\begin{align}
 \int\dd\mu_I(W)\,\dd\widehat\mu(\zeta)
                         \exp\left\{\sum_XV(X)\right\}
       &=\exp\left\{\sum_YE^\#(Y)\right\},
                                      \label{eq:rgnl-log-identity}\\
 \|E^\#(Y)\|_w
       &\le c_S'H_{\mathrm{cl}}\exp(-a_3d_M(Y)).
                                      \label{eq:rgnl-log-bound}
\end{align}
Here $a_3=\kappa-5\kappa_0-4$ and
$c_S'=14c_X^2K_{\mathrm{pol}}^2\exp(3\kappa)$.
The sum of the norms of the terms with $n>N$ in
\eqref{eq:rgnl-log-activity} is at most $2^{-N}$ times the
right side of \eqref{eq:rgnl-log-bound}.
\end{proposition}

\begin{proof}
Set
\[
 \begin{gathered}
 a_4=\kappa-3\kappa_0-3,\quad
 a_5=\kappa-3\kappa_0-2,\quad
 a_6=\kappa-5\kappa_0-3,\\
 H_1=7c_XK_{\mathrm{pol}}\exp(3a_5)H_{\mathrm{cl}},\qquad
 z_S=28H_1\exp(\kappa_0)K_{\mathrm{pol}}\exp(3a_6).
 \end{gathered}
\]
The fixed small-charge choice
\[
 c_S e^{1/4-5\varepsilon}\le1,\qquad
 c_S=10^5c_X^5K_{\mathrm{pol}}^5\exp(6\kappa)
\]
implies
\begin{equation}\label{eq:rgnl-cluster-smallness}
 H_{\mathrm{cl}}\le1,\qquad
 4H_{\mathrm{cl}}K_{\mathrm{pol}}c_X\exp(3a_5)\le1,\qquad
                                      z_S\le\tfrac12.
\end{equation}
Indeed the second left side is
$144c_X^3K_{\mathrm{pol}}^3\exp(3a_5)e^{1/4-5\varepsilon}$.
For the third,
$2z_S=14112c_X^3K_{\mathrm{pol}}^4
\exp(3a_5+3a_6+\kappa_0)e^{1/4-5\varepsilon}$, and
$3a_5+3a_6+\kappa_0\le6\kappa$.

Even elements commute, so expanding their exponential
gives $\exp(\sum_XV(X))=\sum_YK(Y)$.
Moreover $K(Y)$ factors over the connected components
of $Y$: every polymer in its defining set belongs to
exactly one component.
If distinct polymers $X_1,\ldots,X_n$ have connected
union $Y$, their incompatibility graph is connected.
Otherwise two nonempty groups would partition their
union; connectedness of $Y$ would give coincident or
adjacent cubes in different groups and hence an edge
between them.
For each incompatibility-tree edge, choose coincident or
face-adjacent cubes in the two polymers and points
where their nets meet these cubes. The two points
lie in a box of side lengths $M,M,2M$, so they
can be joined by a segment of length at most
$\sqrt6M<3M$. Joining near-minimal nets along
the tree proves
\[
 d_M(Y)\le\sum_i d_M(X_i)+3(n-1).
\]
Using $\|\exp V-1\|\le2\|V\|$ and $a_2=a_5+\kappa_0$
therefore bounds the product by
\[
 \exp(-a_5d_M(Y))
       \big(2H_{\mathrm{cl}}\exp(3a_5)\big)^n
                                  \prod_i \exp(-\kappa_0d_M(X_i)).
\]
Replace distinct sets by ordered tuples with factor $1/n!$,
and use the counting bound for subpolymers of $Y$. The result is
\[
 \|K(Y)\|\le \exp(-a_5d_M(Y))x_Y\exp(x_Y),\qquad
 x_Y=2H_{\mathrm{cl}}K_{\mathrm{pol}}\exp(3a_5)|Y|_M.
\]
By \eqref{eq:rgnl-cluster-smallness},
\[
 \begin{gathered}
 x_Y\le(d_M(Y)+1)/2,\\
 x_Y\le4H_{\mathrm{cl}}K_{\mathrm{pol}}c_X\exp(3a_5)\exp(d_M(Y)/2).
 \end{gathered}
\]
Since $4\exp(1/2)\le7$ and $a_5-1=a_4$, we obtain
\begin{equation}\label{eq:rgnl-Mayer-bound}
 \|K(Y)\|_{(w;\mathsf w_W)}
                         \le H_1\exp(-a_4d_M(Y))=:\mu(Y).
\end{equation}

Weighted Berezin integration costs no norm here:
each paired moment at level zero is
$\pm48^{|J|/2}$, which is exactly covered by the
weight $(\sqrt{48})^{|J|}$ of its generators.
The corresponding moment is $\pm1$ at positive levels.
The photon measure is a probability measure on its compact
box. Hence $\|H(Y)\|_w\le\mu(Y)$.
These integrals produce even elements, since a nonzero
Berezin moment uses an even number of generators.
For compatible polymers the fluctuation modes and
supported photon coordinates are disjoint. The two
product measures therefore factor, with the external
old generators left as coefficients. It follows that
\begin{equation}\label{eq:rgnl-compatible-partition}
 \int\dd\mu_I(W)\,\dd\widehat\mu(\zeta)e^{\sum_XV(X)}
             =\sum_{\{Y_i\}\text{ compatible}}\prod_iH(Y_i).
\end{equation}

We now bound the infinite connected series. Let
\[
 S_n(Y)=\sum_{\bigcup_iY_i=Y}
                   |u^T(Y_1,\ldots,Y_n)|\prod_i\mu(Y_i).
\]
For $n\ge2$, apply \eqref{eq:rgnl-tree-domination}.
For every contributing tree the union estimate and
$a_4=a_6+2\kappa_0$ give
\[
 \prod_i \exp(-a_4d_M(Y_i))
 \le \exp(-a_6d_M(Y))\exp(3a_6(n-1))
                                  \prod_i \exp(-2\kappa_0d_M(Y_i)).
\]
Drop the prescribed union, retaining $Y_i\subseteq Y$,
and root the tree at $1$. Write $q_j$ for the number
of children at $j$. A polymer incompatible with $Y'$
contains one of at most $7|Y'|_M$ cubes consisting of
its cubes and their face neighbors. Since
$|Z|_M\le c_X(d_M(Z)+1)\le\kappa_0(d_M(Z)+1)$ and
$a^q\le q!\exp(a)$, the counting estimate gives
\begin{align*}
 \sum_{Z\text{ incompatible with }Y'}
      \exp(-2\kappa_0d_M(Z))|Z|_M^{q}
            &\le7q!\exp(\kappa_0)K_{\mathrm{pol}}|Y'|_M,\\
 \sum_{Z\subseteq Y}\exp(-2\kappa_0d_M(Z))|Z|_M^q
            &\le q!\exp(\kappa_0)K_{\mathrm{pol}}|Y|_M.
\end{align*}
Summing from leaves to root thus costs at most
$(7\exp(\kappa_0)K_{\mathrm{pol}})^n|Y|_M\prod_jq_j!$.
For prescribed children counts, Cayley's degree formula
\cite[Proposition~4]{PostnikovLin2019}
counts
$(n-2)!/((q_1-1)!\prod_{j\ge2}q_j!)$ trees, where
$q_1\ge1$ and $\sum_jq_j=n-1$.
Consequently
\[
 \sum_{\text{trees}}\prod_jq_j!
   =(n-2)!\sum_{(q_j)}q_1
   \le(n-1)!\binom{2n-2}{n-1}
   \le(n-1)!4^{n-1}.
\]
Collecting the factors, also for $n=1$, gives
\begin{equation}\label{eq:rgnl-connected-n-bound}
 \frac{S_n(Y)}{n!}
 \le7H_1\exp(\kappa_0)K_{\mathrm{pol}}|Y|_M
                                      \exp(-a_6d_M(Y))z_S^{n-1}.
\end{equation}
Since $z_S\le1/2$, sum the geometric series and use
$|Y|_M\le c_X\exp(d_M(Y))$ and $a_6-1=a_3$.
The bound becomes
\[
 98c_X^2K_{\mathrm{pol}}^2\exp(\kappa_0+3a_5)
                                 H_{\mathrm{cl}}\exp(-a_3d_M(Y))
 \le c_S'H_{\mathrm{cl}}\exp(-a_3d_M(Y)),
\]
because $\kappa_0+3a_5=3\kappa-8\kappa_0-6$.
Summing only $n>N$ gives the additional factor $2^{-N}$.
This proves absolute convergence and the claimed estimates.

It remains to justify the logarithm identity after
substitution in \eqref{eq:rgnl-connected-graph-identity}.
The even subalgebra with the coefficient norm is a
commutative Banach algebra, and the polymer set is finite.
The total norm sum of the connected terms is finite by
\eqref{eq:rgnl-connected-n-bound}.
In the exponential of their sum, the norm sum of all
products of connected terms is bounded by the exponential
of that total norm sum. Regrouping is therefore permitted.
Regroup by monomials in the formal variables $z_Y$.
The coefficients are those of
\eqref{eq:rgnl-connected-graph-identity}, namely one
for compatible sets and zero for all other monomials.
Its substitution $z_Y=H(Y)$ gives
\eqref{eq:rgnl-compatible-partition}. Grouping the
connected terms by their union gives
$\sum_YE^\#(Y)$ and proves \eqref{eq:rgnl-log-identity}.
\end{proof}

\subsection{Polynomial approximation at fixed volume}\label{subsec:rgnl-polynomial-approximation}

There are also polynomial approximants to this logarithm
which converge uniformly for all families satisfying the
preceding bound, at fixed side. Replace $\exp V-1$ by
$p_N(V)=\sum_{1\le j\le N}V^j/j!$ in $K,H$, and stop
the connected series at $n=N$; denote the result $E_N^\#$.
It is a polynomial in integrals of products of $V(X_i)$,
$X_i\subseteq Y$, and has the same bound
\eqref{eq:rgnl-log-bound}. To verify uniform convergence,
write $D=|\mathcal P_k|$. Since $\|V(X)\|\le1$,
\[
 \|\exp V(X)-1-p_N(V(X))\|
                                  \le\frac{\exp(1)}{(N+1)!}.
\]
Telescoping each finite product in $K$ gives
\[
 \|H(Y)-H_N(Y)\|\le
                             \frac{4^DD\exp(1)}{(N+1)!}.
\]
Thus this difference is at most $\delta_N\mu(Y)$, where
\[
 \delta_N=\frac{4^DD\exp(1)}
                      {(N+1)!\min_{Y\in\mathcal P_k}\mu(Y)}
                                                   \longrightarrow0.
\]
Replacing the $n$ factors in a connected term costs at
most $n\delta_N\prod_i\mu(Y_i)$. Summing
\eqref{eq:rgnl-connected-n-bound} with this factor, using
$\sum_{n\ge1}n2^{-(n-1)}=4$, and adding the tail yields
\[
 \sup_Y\|E^\#(Y)-E_N^\#(Y)\|
 \le28\delta_NH_1\exp(\kappa_0)K_{\mathrm{pol}}(\ell/M)^3
                                  +2^{-N}c_S'H_{\mathrm{cl}}.
\]
This convergence is uniform in the connections and the
allowed data. The displayed error depends on the side;
no volume-independent approximation is claimed.

\section{Symmetries of the local logarithm}
\label{subsec:rgnl-symmetries}

Apply \cref{prop:rgnl-cluster-logarithm} with
$V=F$, $w=h/4$ and denote its output by
$E_k^\#(m,\mathscr E)$. Each connected term assigned to
$Y$ is a product of integrated activities on subpolymers
of $Y$. Thus ordinary and admissible pair support, local
connection dependence, and dependence only on $m'$ and
$\mathscr E^r(X)$, $X\subseteq Y$, pass to the sum.
The uniform bound and the polynomial approximants prove
analyticity of its coefficients in $(A,\mathscr C)$.

Let $\mathbb G_k^L$ denote the chamber symmetry group with
translations in $LM\Z^3$, so that it preserves every
parameter partition. Its actions and the gauge,
charge-conjugation and reality actions are those of
\cref{subsec:rg-activity-spaces}.

\begin{proposition}\label{prop:rgnl-symmetries}
The local logarithm satisfies
\begin{align}
 a_{\tilde g}E_k^\#(\tilde gY;\tilde gA,\tilde g\mathscr C)
          &=E_k^\#(Y;A,\mathscr C),\qquad \tilde g\in\mathbb G_k^L,
                                      \label{eq:rgnl-lattice-symmetry}\\
 G_{\lambda+\lambda'}E_k^\#(Y;A^1,\mathscr C^1)
          &=E_k^\#(Y;A,\mathscr C),    \label{eq:rgnl-gauge-symmetry}\\
 \mathrm{Cg}\,E_k^\#(Y;-A,-\mathscr C)
          &=E_k^\#(Y;A,\mathscr C).    \label{eq:rgnl-charge-symmetry}
\end{align}
In the gauge identity both pairs belong to the domain and
 $A^1=A-d\lambda$, $\mathscr C^1=\mathscr C-d\lambda'$
on all bonds at distance at most $1$ from $Y$, with
real $\lambda,\lambda'$. Here $d\lambda$ denotes the
lattice gradient, as in the preceding gauge conventions.
For real $m$ and $\mathscr E\in K_k^x$ one also has
\begin{equation}\label{eq:rgnl-reality-symmetry}
 \mathbf K E_k^\#(Y;-\overline A,-\overline{\mathscr C})
                                        =E_k^\#(Y;A,\mathscr C).
\end{equation}
Every monomial is balanced.
\end{proposition}

\begin{proof}
We first record how transformations pass through the
finite operations. Each of the four algebra maps is a
continuous ring automorphism, linear for the first three
and antilinear for $\mathbf K$. If an odd field transforms
like the generators it replaces, the two compositions of
the field substitution with the automorphism agree on
each generator; hence they agree on the algebra.
The same observation applies to adding a field.
Real-coefficient finite differences
\eqref{eq:rgnl-finite-difference} then commute with these
actions. A lattice symmetry also permutes the parameter
cubes and localization summands. It preserves face
connectedness, incidence relations and the Ursell coefficient.

The contour functional commutes with the linear actions.
For the antilinear action, parametrize $t=r\exp(i\theta)$:
\[
 \mathcal C_rf=\frac1{2\pi}\int_0^{2\pi}
             \frac{f(r\exp(i\theta))}{r\exp(i\theta)-1}\,\dd\theta.
\]
Conjugation of coefficients, together with
$t\mapsto\bar t$ and $\theta\mapsto-\theta$, gives the
same integral at the transformed other arguments.
The integration over $W$ commutes with every linear or
antilinear action which sends
$W\mapsto UW$, $\bar W\mapsto(U^{-1})^T\bar W$.
Indeed Wick's rule preserves every paired moment under
this change, and all these moments are real.
This applies to the lattice and gauge actions, to
$\mathbf K$ with $U=\varkappa$, and to charge conjugation
extended by the identity on $W$.
The last extension will require a separate covariance
argument below.

For a lattice symmetry the primitive covariance gives
\[
 Z(\tilde gs')[\tilde g\zeta]=\tilde gZ(s')[\zeta],\qquad
 H_k(\tilde gs'',\tilde gB)\tilde g=\tilde gH_k(s'',B),
 \qquad
 \Gamma_k(\tilde gs'',\tilde gB)\tilde g=\tilde g\Gamma_k(s'',B).
\]
The barred action is $(S^{-1})^T$ on spinors;
$(S^{-1})^T=\pm\mathsf C^{-1}S\mathsf C$ gives the
same covariance for $\bar H_k$.
H\"older quotients commute with the simultaneous
transformation of their endpoints and connection.
Thus $J_k^\#\Psi$ and $W_k^\#$ transform like their
ordinary and pair generators.
The quadratic primitive satisfies
$M_k(\tilde gX,\tilde gB)=\tilde gM_k(X,B)\tilde g^{-1}$,
so its bilinear form is invariant.
The mass form is invariant by the change of chamber
site and spin indices, and $R\mathscr E$ has the
incoming lattice symmetry. The reconstruction substitution
commutes with the action by
\eqref{eq:rg-step-substitution-covariances}.
The box measure is invariant under $\zeta\mapsto\tilde g\zeta$.
Apply these identities to $G_1,G_2,G_3$, the first
localization and substitution, then to the second
localization, $K,H$, and the convergent connected sum.
They prove \eqref{eq:rgnl-lattice-symmetry}.

For the gauge identity set $\Lambda=\lambda+\lambda'$.
In a summand assigned to $Y$, all parameters outside $Y$
vanish. The local row result proves that $B^1$ may be
replaced by $B-d\Lambda$ in every row used in that
summand. The latter field is still in the operator domain,
since real gauge shifts preserve it.
Telescoping the gradient along each path gives
\[
 U_{B-d\Lambda}(x,x')
       =U_B(x,x')\exp\big(-ie(\Lambda(c')-\Lambda(c))\big).
\]
The ordinary primitive row obeys
\[
 H_k(s'',B^1)(x,\cdot)
       =\exp(ie\Lambda(c))H_k(s'',B)(x,\cdot)
                                                 \exp(-ie\Lambda^u),
\]
when it is nonzero, where $\Lambda^u$ is the function
on the unit indices. The barred row has the inverse
phases. In the H\"older row the phase at its second
point cancels that from the transport, leaving exactly
$\exp(ie\Lambda(c))$, the phase of the first point.
The row of $\Gamma_k$ has the corresponding two-sided
gauge covariance. Its composition with $H_k^\#$ thus
gives the same covariance for $W_k^\#$, including the
rows which vanish outside $Y$.
The primitive bilinear and the local mass form are
gauge invariant, and so is $R\mathscr E$ under the
separate shifts of its two slots. The reconstruction
substitution also commutes with this local gauge action.
The photon fields and photon measure are unchanged.
The finite operations and integration therefore prove
\eqref{eq:rgnl-gauge-symmetry}.

For reality, the photon primitives are real on real
parameter families. Their analytic dependence extends
this identity to
$\overline{O(\bar s')}=O(s')$ on the full polydisc.
Since $C$ is real,
\[
 Z(\bar s')[-\zeta]=-\overline{Z(s')[\zeta]}.
\]
Transform $s\mapsto\bar s$, $t\mapsto\bar t$,
$\zeta\mapsto-\zeta$ and
$(A,\mathscr C)\mapsto(-\bar A,-\overline{\mathscr C})$.
The fermion identities are
\[
 \overline{H_k(\bar s'',-\bar B)}
                         =\varkappa^{-1}H_k(s'',B)\varkappa,\qquad
 \overline{\Gamma_k(\bar s'',-\bar B)}
                         =\varkappa^{-1}\Gamma_k(s'',B)\varkappa.
\]
For the barred field use
$\varkappa\mathsf C=-\mathsf C\varkappa$.
The transport relation
$\overline{U_{-\bar B}(x,x')}=U_B(x,x')$
then gives the same identities for the H\"older rows.
The local quadratic primitive obeys the corresponding
identity, and $\varkappa^T\varkappa=1$ preserves the
mass form. The number $m'$ is real for the present data,
and $R\mathscr E$ and the reconstruction substitution
have the reality symmetry. The contour argument just
given, invariance of the box under sign reversal, and
weighted integration prove
\eqref{eq:rgnl-reality-symmetry}.

Charge conjugation gives the first-localization symmetry
by
\[
 \begin{gathered}
 H_k(s'',-B)\mathsf C=\mathsf C\bar H_k(s'',B),\qquad
 \bar H_k(s'',-B)\mathsf C=\mathsf C H_k(s'',B),\\
 M_k(X,-B)^T=\mathsf C^{-1}M_k(X,B)\mathsf C .
 \end{gathered}
\]
Anticommuting the two transformed legs verifies the
bilinear and mass-form identities. The first-step
vertex has this symmetry by its explicit two-half formula.
It follows that
\[
 \mathrm{Cg}\,E_{\le3}(X;-A,-\mathscr C,-\zeta)
                                        =E_{\le3}(X;A,\mathscr C,\zeta).
\]
Direct covariance of the second substitution is not
available, since its unbarred and barred fluctuations
use different factors. We instead compare their Wick
covariances.

Introduce independent generators $\phi^\#$ indexed like
$\psi^\#$, with their own charge-conjugation action,
and let $\operatorname{ad}$ add them to the old generators.
Put
\begin{align*}
 \mathcal V(X)
   ={}&\operatorname{ad}E_{\le3}(X)
      +\operatorname{ad}\mathscr E^r(X)-\mathscr E^r(X)\\
     &-\mathbf1_{X=\square}m'
          \big(\operatorname{ad}
                         \langle\bar\psi,\psi\rangle_\square
                                -\langle\bar\psi,\psi\rangle_\square\big).
\end{align*}
It has charge-conjugation symmetry under
$(A,\mathscr C,\zeta)\mapsto(-A,-\mathscr C,-\zeta)$.
Let $I_{s'',B}$ fix the old generators and replace
$\phi^\#$ by $W_k^\#(s'',B)$. The second-localization
input is exactly $I_{s'',B}\mathcal V(X)$, since
the polynomial contours in $G_4,G_5$ give their value
at $1$ minus their value at $0$.
For any union $U$, interpolation restricted to cubes
of $U$, with the disconnected terms removed, gives
\begin{equation}\label{eq:rgnl-subvolume-substitution}
 \sum_{Y\subseteq U}F(Y)=
                        I_{\mathbf1_U,B}\sum_{X\subseteq U}\mathcal V(X).
\end{equation}

At $s''=\mathbf1_U$ define
\[
 \phi_\pm=H_k(\mathbf1_U,\pm B)\Gamma_k(\mathbf1_U,\pm B)W,
 \qquad
 \bar\phi_\pm=\bar H_k(\mathbf1_U,\pm B)\bar W,
 \qquad
 \phi'=\mathsf C\bar\phi_-,\quad
 \bar\phi'=\mathsf C^{-1}\phi_-.
\]
The substitution by $(\phi',\bar\phi')$ and their
H\"older quotients, denoted $I'$, satisfies
$\mathrm{Cg}\,I_{\mathbf1_U,-B}=I'\mathrm{Cg}$
on every generator, when $\mathrm{Cg}$ fixes $W$.
The only nonzero ordinary pair expectations for the
positive substitution are the kernel of
\[
 \mathcal G_U(B)
        =H_k(\mathbf1_U,B)\Gamma_k(\mathbf1_U,B)
                                            \bar H_k(\mathbf1_U,B)^T.
\]
The identity moment has kernel
$\mathbf1_{u=u'}/\nu(u)$, with $\nu=1/48$ at zero.
This precisely cancels one of the two measures in
the kernel composition, so the same formula holds
at every level. The primitive transpose relation gives
\[
 \mathcal G_U(-B)^T
                      =-\mathsf C^{-1}\mathcal G_U(B)\mathsf C^{-1}.
\]
Anticommuting $\bar\phi_-$ and $\phi_-$ therefore gives
\[
 \int\dd\mu_I(W)\,\phi'(x)\bar\phi'(x')^T
     =\mathsf C[-\mathcal G_U(-B)^T(x,x')]\mathsf C
     =\mathcal G_U(B)(x,x').
\]
Same-bar pair expectations vanish in both substitutions.
The pair generators are obtained by the identical
linear H\"older maps at $B,-B$, so their pair expectations
also agree. Wick's theorem now proves
$\int\dd\mu_I(W)I'Q=\int\dd\mu_I(W)I_{\mathbf1_U,B}Q$
for every polynomial $Q$ in the old and $\phi^\#$
generators: each monomial integral is the Pfaffian of
its pair expectations, with old generators as coefficients.

Apply this equality to the exponential of
$\sum_{X\subseteq U}\mathcal V(X)$ and use
$\zeta\mapsto-\zeta$. By
\eqref{eq:rgnl-subvolume-substitution}, the integrated
subvolume expression
$\Xi_U=\int\dd\mu_I(W)\dd\widehat\mu(\zeta)
\exp(\sum_{Y\subseteq U}F(Y))$
has charge-conjugation symmetry.
But $\Xi_U=\sum_{Z\subseteq U}H(Z)$, by the finite
Mayer expansion. Inverting this finite subset sum gives
\[
 H(Z)=\sum_{U\subseteq Z}(-1)^{|Z\setminus U|_M}\Xi_U.
\]
Thus each $H(Z)$ has the symmetry, and the real
coefficients and absolute convergence of the connected
series prove \eqref{eq:rgnl-charge-symmetry}.

Finally take a constant $\lambda$ and $\lambda'=0$
in \eqref{eq:rgnl-gauge-symmetry}. A monomial with
$n_u$ unbarred and $n_b$ barred legs acquires the factor
$\exp(ie\lambda(n_u-n_b))$. Invariance for every real
$\lambda$ and independence of the basis monomials force
its coefficient to vanish unless $n_u=n_b$.
\end{proof}

\section{Scaling and coupling estimates}
\label{subsec:rgnl-scaled-correction}

Let $S_L$ be the field scaling of
\exteq{30}{NormalizedContraction}. The nonlinear correction is
\begin{equation}\label{eq:rgnl-scaled-correction}
 E_k^*(m,\mathscr E)(Y;\mathscr A,\mathscr C)
  =\sum_{\bar X^L=LY}S_L
        E_k^\#(m,\mathscr E)(X;\mathscr A_L,\mathscr C_L).
\end{equation}
This formula uses field scaling and summation over fine
polymers. Quadratic adjustment enters the incoming
activity term in the full step below.
The domain inclusions
\exteq{29}{NormalizedContraction} ensure
$(\mathscr A_L,\mathscr C_L)
\in\widetilde R_k^-\times\mathfrak X_k$.

\begin{proposition}\label{prop:rgnl-scaled-correction}
For $(m,\mathscr E)\in\operatorname{Ball}_k^x(4)$,
\[
 E_k^*(m,\mathscr E)\in K_{k+1}^{x,\C},\qquad
                  \|E_k^*(m,\mathscr E)\|_{k+1}^x
                                      \le e^{1/4-6\varepsilon}.
\]
For real data its value belongs to $K_{k+1}^x$.
On any affine complex line in the data space its value
is norm continuous on the part in the closed ball and
norm analytic in its interior.
For two data in $\operatorname{Ball}_k^x(1)$,
\begin{equation}\label{eq:rgnl-coupling-Lipschitz}
 \|E_k^*(m_1,\mathscr E_1)-E_k^*(m_2,\mathscr E_2)\|_{k+1}^x
 \le e^{1/4-6\varepsilon}
       \left(\frac{|m_1-m_2|}{w_k^{\mathrm m}}
              +\frac{\|\mathscr E_1-\mathscr E_2\|_k^x}
                                                    {w_k^{\mathrm E}}\right).
\end{equation}
The second quotient is absent at $k=0$.
Restriction to $\mathscr C=0$ gives the correction formed
without an external slot from $(m,\mathscr E|_0)$.
\end{proposition}

\begin{proof}
The scaling identity gives
\[
 \|S_LQ\|_{h_{k+1}}
   =\|Q\|_{(L^{-1}h_{k+1,1},L^{-1-\alpha}h_{k+1,2})}
   \le\|Q\|_{h_k/4},
\]
since $h_{k+1,1}=L^{-1/8}h_{k,1}$,
$h_{k+1,2}=L^{-1/8+\varepsilon/2}h_{k,2}$ and $L\ge4$.
The polymer sum in \eqref{eq:rgnl-scaled-correction}
uses $Ld_M(Y)\le d_M(X)$ and $X\subseteq LY$.
Splitting $a_3=(a_3-\kappa_0)+\kappa_0$ therefore yields
\begin{align*}
 \sum_{\bar X^L=LY}\exp(-a_3d_M(X))
  &\le \exp(-L(a_3-\kappa_0)d_M(Y))
                         \sum_{X\subseteq LY}\exp(-\kappa_0d_M(X))\\
  &\le c_XK_{\mathrm{pol}}L^3
                  \exp(-[L(\kappa-6\kappa_0-4)-1]d_M(Y))\\
  &\le c_XK_{\mathrm{pol}}L^3\exp(-\kappa d_M(Y)).
\end{align*}
Together with \eqref{eq:rgnl-log-bound} this gives
\[
 \|E_k^*\|_{k+1}^x
  \le c_XK_{\mathrm{pol}}L^3c_S'H_{\mathrm{cl}}
  =504c_X^5K_{\mathrm{pol}}^5\exp(3\kappa)L^3e^{1/4-5\varepsilon}
  \le e^{1/4-6\varepsilon},
\]
using $c_SL^3e^\varepsilon\le1$.

All generators of a summand lie over $X$ or at a pair
admissible for $X$. Scaling maps the ordinary sites
into $Y$; it maps each of the six paths into the
corresponding scaled path and maps
$(x,x^{(3)})$ to $(x/L,(x/L)^{(3)})$.
These pairs are admissible for $Y$, including
intracluster pairs. A bond meeting $X$ scales to
one meeting $Y$, giving local connection dependence.
The scaled-domain inclusion and coefficient analyticity
give analytic kernels.
The map $(S,a)\mapsto(S,a/L)$ identifies
$\mathbb G_k^L$ with $\mathbb G_{k+1}$ and intertwines
its generator actions with $S_L$. For gauge transformations
use $\Lambda_L(c)=L^{1/2}\Lambda(c/L)$ and
$(d\Lambda)_L=d(\Lambda_L)$.
A gauge relation on the unit neighborhood of $Y$
then holds on the unit neighborhood of every
$X\subseteq LY$. Charge conjugation and reality also
commute with the real scaling factors.
\Cref{prop:rgnl-symmetries} thus gives all the
defining symmetries of $K_{k+1}^{x,\C}$, and the real
space assertion follows from its reality clause.

For norm analyticity, take data
$(m_a+t m_b,\mathscr E_a+t\mathscr E_b)$ in the closed
ball and denote their parameter set by $T$.
It is closed and convex. The finite activity $F_t(X)$
is affine in $t$ by \cref{prop:rgnl-finite-bounds}.
The polynomial approximants to $E_k^\#$ are therefore
polynomials in $t$. Reblock them by
\eqref{eq:rgnl-scaled-correction}; the resulting polynomials
$f_N(t)$ converge uniformly on $T$ in the full activity
norm, since both finite polymer sets are fixed and the
error bound in \cref{subsec:rgnl-polynomial-approximation} is uniform.
Their values have bounded full norm.
If $T$ has more than one point, it contains infinitely
many points. The coefficients of each polynomial are
finite numerical linear combinations of its values at
distinct points of $T$, by polynomial interpolation.
Hence these coefficients have finite full norm, and the
uniform limit is norm continuous on $T$.
If $T$ is empty or a singleton there is no continuity
or interior analyticity assertion to check.

For a closed disk $|t-t_0|\le r$ in $T$, take its
coefficientwise Cauchy integrals
\[
 f_j=\frac1{2\pi i}\int_{|t-t_0|=r}
                             f(t)(t-t_0)^{-j-1}\,\dd t.
\]
They have analytic connection kernels, and the support,
symmetry and balance conditions pass through this
complex linear integral. Moreover
$\|f_j\|_{k+1}^x\le e^{1/4-6\varepsilon}r^{-j}$.
Coefficientwise Cauchy's formula, with this summable
norm bound, gives
$f(t)=\sum_{j\ge0}f_j(t-t_0)^j$ in activity norm for
$|t-t_0|<r$. This proves the stated analyticity.

To obtain \eqref{eq:rgnl-coupling-Lipschitz}, vary first
the mass, keeping $\mathscr E_1$ fixed. With
$\Delta m=m_1-m_2\ne0$, the line
$(m_2+t\Delta m,\mathscr E_1)$ remains in the large
ball for $|t|\le r=3w_k^{\mathrm m}/|\Delta m|$.
Here $r\ge3/2$. The contour formula on any circle
$1<r'<r$ gives
\[
 \|f(1)-f(0)\|_{k+1}^x
        \le\frac{e^{1/4-6\varepsilon}}{r'-1}.
\]
Let $r'\uparrow r$ and use $r-1\ge r/3$.
The bound is the first quotient on the right of
\eqref{eq:rgnl-coupling-Lipschitz}.
Vary the activity with mass $m_2$ fixed and radius
$3w_k^{\mathrm E}/\|\mathscr E_1-\mathscr E_2\|_k^x$
to obtain the second quotient.
The zero differences require no contour.

Finally restriction commutes with extraction:
$m(\mathscr E)=m(\mathscr E|_0)$ and
$(R\mathscr E)|_0=R(\mathscr E|_0)$.
Every fermion operator uses the sum of its first
connection and $\mathscr C$; photon operations and
measures do not use $\mathscr C$.
At $\mathscr C=0$ the two constructions therefore have
identical finite expressions, connected series and
scaled corrections. Analyticity without a slot is
understood in the first connection alone.
\end{proof}

The possible H\"older legs can also be read directly.
In $G_2,G_4$ an existing pair generator is changed by
addition of a fluctuation field; this creates no new
pair index. The mass contours contain no pair
generators. The reconstruction substitution and the
first-step vertex introduce only pairs $(x,x^{(3)})$.
Every later sum, product and integral preserves this
set of available indices. Thus the correction has
only scaled incoming pair indices and these scaled
type-$3$ indices. At $k=0$ it has only the latter,
of length $\eta_1$, since the incoming activity is zero.

\section{Identification with the exact density}
\label{subsec:rgnl-current-step}

The three terms in \eqref{eq:rgnl-full-step} have different origins:
the normalized incoming activity, the local determinant difference,
and the fluctuation logarithm. We first bound the determinant term,
then identify their evaluated sum with the exact block density.

The reblocked determinant term belongs to $K_{k+1}^x$.
Indeed the scaled arguments lie in its inverse domain,
and each summand is analytic and reads only fine
bonds meeting $X$. These scale to bonds meeting $Y$.
There are no legs. The lattice symmetry permutes the
reblocking summands; local gauge shifts of the two
slots become the corresponding real shifts on these
fine read bonds, where determinant locality and gauge
invariance apply. Evenness of $E_k^d$ gives charge
conjugation, and its conjugation identity gives reality.
The pointwise bound
\exteq{82}{FermionBlockManuscript} and
\exteq{83}{NormalizedContraction} show
\begin{align}
 \|\mathcal L\delta E_k^{\mathrm{det}}\|_{k+1}^x
    &\le8c_{\mathrm{rb}}L^3C_dM^3e_k^{1/4}
      \le e_k^{1/4-2\varepsilon}
      \le e_{k+1}^{1/4-2\varepsilon}.
                                      \label{eq:rgnl-reblocked-determinant}
\end{align}
Here $L(\kappa-\kappa_0-1)\ge\kappa$ controls the
remaining exponential weight, and the fixed charge
condition is
$8c_{\mathrm{rb}}L^3C_dM^3e_k^{2\varepsilon}\le1$.
Restriction to $\mathscr C=0$ gives the same result
for $\mathcal LE_k^{\mathrm{det}}$.

\begin{proof}[Proof of \cref{thm:rgnl-current-step}]
Use the notation of \cref{prop:rg-step-density-identity},
with $A=\mathscr A'_L$, $B=A+\mathscr C$ and
$Z=Z(1)$ the full-strength photon fluctuation.
The primitive families satisfy their exact identities
at full strength:
$H_k(1,B)=H_k(B)$, $\Gamma_k(1,B)=\Gamma_k(B)$, and
$\sum_XM_k(X,B)=b_k-D_k(B)$ for $k\ge1$.
The induced fields obey
\[
 \psi^{0,\#}=H_k^\#(B)\Psi^c,\qquad
 H_k^\#(B)\Psi^f=\psi^{0,\#}+W_k^\#(1,B).
\]
Define three evaluation homomorphisms. Let $\operatorname{ev}$
send old generators to $\psi^{0,\#}$ and fix $W$;
let $\operatorname{ev}_1$ send them to $H_k^\#(B)\Psi^f$;
for $k\ge1$, let $\operatorname{ev}_f$ additionally send
$\Psi,\bar\Psi$ to $\Psi^f,\bar\Psi^f$.
At zero take $\operatorname{ev}_f=\operatorname{ev}_1$.
These maps exist since their images are odd and have
the required antiperiodicity. They commute with the
finite coefficient integrals; $\operatorname{ev}$ also
commutes with Berezin integration, since it fixes $W$.

On the first contours $\mathcal C_rf=f(1)-f(0)$.
For $G_1$, summing the local primitives gives the
Dirac difference
$-\langle\bar\Psi,(D_k(B+Z)-D_k(B))\Psi\rangle_1$.
For $G_2,G_3$, evaluating
$\psi^\#+J_k^\#(1,B,Z)\Psi$ gives
$H_k^\#(B+Z)\Psi^f$.
The evaluations of their contour differences therefore give
\begin{equation}\label{eq:rgnl-first-contour-evaluation}
 \operatorname{ev}_f\sum_X\sum_{i\le3}G_i(1,X)
  =\mathrm{Act}_k^x[m',\mathscr E^r](A+Z,\mathscr C;\Psi^f)
       -\mathrm{Act}_k^x[m',\mathscr E^r](A,\mathscr C;\Psi^f).
\end{equation}
At zero the explicit first vertex gives this same
identity with the chamber pairing: $D_0(B)=D_B+\mu$
and the mass form has no connection dependence.

The reconstruction left inverse gives, on every local
first expression,
$\operatorname{ev}_1\operatorname{sub}_{Y,B}
=\operatorname{ev}_f$.
This follows on the $\Psi$ generators from $T_k(B)H_k(B)=1$
and on the old generators from their definition.
Adding $W_k^\#(1,B)$ to $\psi^\#$ in a polynomial and then
applying $\operatorname{ev}$ gives
$\operatorname{ev}_1$ applied to that polynomial.
Apply these identities to
\eqref{eq:rgnl-two-localization-sums}.
The evaluated $G_{\le3}^{\,W}$ sum is
\eqref{eq:rgnl-first-contour-evaluation}.
The evaluated polynomial contours $G_4,G_5$ give
\[
 I_k^x[m',\mathscr E^r](A,\mathscr C;\Psi^f)
                      -I_k^x[m',\mathscr E^r](A,\mathscr C;\Psi^c).
\]
Consequently
\[
 \operatorname{ev}\sum_YF(Y)=V_k^x,\qquad
 \Xi_k^x
     =\exp\left\{\sum_X
                  \operatorname{ev}E_k^\#(m,\mathscr E)(X;A,\mathscr C)\right\}.
\]
The second equality is
\eqref{eq:rgnl-log-identity} evaluated by
$\operatorname{ev}$; its evaluation commutes with the
integrals as just noted.

The other remaining factor is the determinant quotient.
Equation~\exteq{78}{FermionBlockManuscript} gives
\[
 \frac{\delta Z_k(B)}{\delta Z_k(0)}
      =\exp\left\{\sum_X\delta E_k^{\mathrm{det}}(X;A,\mathscr C)
                               +\sum_XE_k^{\mathrm{det}}(X,\mathscr C)\right\}.
\]
On inserted fields adjustment changes no value:
$\operatorname{ev}R\mathscr E(X)
=\operatorname{ev}(R\mathscr E)^{\mathrm{nat}}(X)$.
The equality is the telescoping reconstruction identity
for a small polymer and the definition for a large one.
Finally $\operatorname{ev}=\operatorname{ev}'S_L$,
where $\operatorname{ev}'$ inserts the outgoing fields.
Each fine polymer has exactly one outgoing reblocking
polymer. Thus the three activity sums combine into the
outgoing $\mathscr E^+$, and the two scalar external sums
combine into \eqref{eq:rgnl-external-scalar-map}.
Insert these identities in
\eqref{eq:rg-step-density-identity}. All exponents are
even and commute, so their sums give
\eqref{eq:rgnl-exact-step}, with exactly its previously
derived mass, vacuum and determinant normalization.

Integrate over the outgoing generators and apply the exact
marginal \eqref{eq:rg-step-block-marginal}. This gives
\eqref{eq:rgnl-integrated-step}.
The integrand is admissible throughout the photon box:
Theorem~\extref{B}{MaxwellBlockManuscript} places
$A_k^{\min}(A')+\mathrm{Sq}_k\zeta$ in
$\mathrm{Dom}_k$, and its reconstructed connection is
an affine, continuous function of $\zeta$.
The density coefficients are analytic in this connection
on its operator and activity domains, hence continuous
and integrable on the compact box.
The external scalar factor is fixed during this integral.
No condition on the harmonic part of $A'$ enters.

The three terms of \eqref{eq:rgnl-full-step} belong to
$K_{k+1}^x$ by Proposition~\extref{6.1}{NormalizedContraction},
\eqref{eq:rgnl-reblocked-determinant}, and
\cref{prop:rgnl-scaled-correction}.
Their respective norms are bounded by the three terms
in \eqref{eq:rgnl-activity-step-bound}.
Equation~\exteq{46}{NormalizedContraction} bounds the
first term with the stronger exponent $3\varepsilon/2$.
The mass and vacuum extraction bounds give
\eqref{eq:rgnl-mass-step-bound} and
\eqref{eq:rgnl-vacuum-step-bound}, since
\[
 0\le\delta\varepsilon_k^0
        \le12\exp(-p_{0,k}^2/2)\le e_k^7.
\]
At zero all incoming extracted terms vanish.
The exact algebra uses all $192n_0$ pairs with the
weight $1/48$ and integrates to the $4n_1$-pair
outgoing density. There is no removed multiplicity
or additional vacuum normalization. The displayed
zero-level formulas and bounds follow.
\end{proof}

The restriction to zero external connection commutes
with all three terms of \eqref{eq:rgnl-full-step}.
The determinant difference there becomes
$E_k^d(X,\mathscr A)-E_k^d(X,0)$, and the external scalar
map preserves $\mathfrak P_k(\cdot,0)=0$.
Thus \eqref{eq:rgnl-exact-step} and
\eqref{eq:rgnl-integrated-step} also give the step without
a slot, with incoming activity $\mathscr E|_0$.

The displayed norm estimates depend on the running charge
and the fixed constants, and on no side or harmonic part.
They were proved for any primitive system with the stated hypotheses on
the present torus. Identifying the outputs from different
systems or different volumes requires the separate
comparison arguments; it does not follow from these bounds.


\printbibliography
\end{document}
